
% ============================================================
% Appendix: Theoretical Guarantees for Contextual Hamiltonian Enrichment
% Drop this file into the NeurIPS source and include with:
%   \input{appendix_theory}
%
% Required packages in the main preamble:
%   \usepackage{amsmath,amssymb,amsthm,mathtools}
%   \usepackage{natbib}
%
% Optional theorem environments if not already defined:
%   \newtheorem{theorem}{Theorem}
%   \newtheorem{lemma}{Lemma}
%   \newtheorem{proposition}{Proposition}
%   \newtheorem{corollary}{Corollary}
%   \theoremstyle{definition}
%   \newtheorem{assumption}{Assumption}
%   \newtheorem{remark}{Remark}
% ============================================================

\section{Theoretical Guarantees for Contextual Hamiltonian Enrichment}
\label{app:theory}

This appendix collects the formal properties used in Section~3.6.
The goal is not to claim unconditional safety or global optimality. Instead, we prove
that contextual Hamiltonian enrichment has five conditional properties aligned with
the central claims of the paper: (i) an added energy term induces a new force and
sensitivity channel; (ii) the enriched system preserves the port-Hamiltonian
dissipativity structure; (iii) the discrete rollout is practically dissipative for small
enough step size; (iv) when context gradients vanish, the context-enriched field remains close to the
geometry-only policy and cannot generate a spurious lateral escape; and (v) when a
projected risk-gradient margin exists, the contextual force activates and locally
reduces risk. We also justify the detached-quantile CVaR gradient used in the
episode loop and state the constrained risk-averse OCP interpretation used by the
upper-confidence update. The passivity arguments follow the standard port-Hamiltonian energy
balance~\citep{vanderschaft2014port,desai2021porthamiltonian}; the CVaR argument
uses the Rockafellar--Uryasev variational representation~\citep{rockafellar2000optimization}
and standard empirical quantile sensitivity arguments~\citep{hong2014monte}.

\subsection{Setup and notation}
\label{app:setup}

Let \(q \in \mathbb R^d\) denote the configuration coordinate and \(p\in\mathbb R^d\)
its cotangent variable. Let \(x=(q,p)\), let \(M\succ 0\) be the mass matrix, and
write
\[
    v = M^{-1}p .
\]
The geometry-only Hamiltonian is
\[
    H_{\rm scaf}(q,p)
    =
    \frac12 p^\top M^{-1}p + H_{\mathrm{geom}}(q),
\]
where \(H_{\mathrm{geom}}\) contains the goal and geometric barrier terms. The context-enriched field
adds the context Hamiltonian
\[
    H_{\mathrm{ctx}}(q;z)
    =
    \lambda_s \widetilde r(\mathbf c;z)
    +
    \lambda_h b_{\mathrm{sp}}(\phi(\mathbf c;z)),
\]
where \(\mathbf c\) is the workspace position component of \(q\), \(z\) denotes the
locally sensed context patch, \(\widetilde r\) is the smoothed soft-risk field, and
\(\phi\) is the signed-distance field for hard hazards. Thus
\[
    H_{\rm enr}(q,p;z)
    =
    H_{\rm scaf}(q,p)+H_{\mathrm{ctx}}(q;z).
\]
The corresponding context force is
\[
    F_{\mathrm{ctx}}(q;z)
    =
    -\nabla_q H_{\mathrm{ctx}}(q;z)
    =
    -\lambda_s \nabla_q \widetilde r(\mathbf c;z)
    -
    \lambda_h b_{\mathrm{sp}}'(\phi(\mathbf c;z))\nabla_q \phi(\mathbf c;z).
\]
This is exactly the soft-risk plus hard-hazard force decomposition used in the main
paper.

The continuous-time enriched port-Hamiltonian dynamics are
\begin{align}
    \dot q &= \nabla_p H_{\rm enr}(q,p;z)=M^{-1}p, \label{eq:ct_q}\\
    \dot p &=
    -\nabla_q H_{\rm enr}(q,p;z)
    -
    D(q,p;z)M^{-1}p
    +
    B(q,p;z)u , \label{eq:ct_p}
\end{align}
where \(D(q,p;z)\succeq 0\) is the dissipative port. The rollout used in the paper is
the semi-implicit update
\begin{align}
    p_{t+1}
    &=
    p_t
    -\tau\nabla_q H_{\mathrm{geom}}(q_t)
    -\tau\nabla_q H_{\mathrm{ctx}}(q_t;z_t)
    -\tau D_tM^{-1}p_t
    +\tau B_tu_t, \label{eq:disc_p}\\
    q_{t+1}
    &=
    q_t+\tau M^{-1}p_{t+1}. \label{eq:disc_q}
\end{align}
When \(D_tM^{-1}p_t\) is implemented as the scalar dissipative port
\(R_\theta(q_t,p_t)=\gamma M^{-1}p_t\), this reduces to the update in the main text.

\begin{assumption}[Rollout regularity]
\label{ass:regularity}
There exists a compact rollout tube \(\mathcal K\subset \mathbb R^{2d}\) such that
the geometry-only and context-enriched trajectories remain in \(\mathcal K\) for the horizon under
analysis. On \(\mathcal K\), \(H_{\mathrm{geom}}\) and \(H_{\mathrm{ctx}}\) are \(C^2\),
and their gradients are Lipschitz.
\end{assumption}

\begin{assumption}[Bounded context force]
\label{ass:bounded_ctx}
There exists \(G_{\mathrm{ctx}}<\infty\) such that
\[
    \|\nabla_q H_{\mathrm{ctx}}(q;z)\|\le G_{\mathrm{ctx}}
\]
for all \((q,p)\in\mathcal K\) and all context patches considered.
\end{assumption}

\begin{assumption}[Dissipative port]
\label{ass:diss}
There exists \(\rho\ge 0\) such that
\[
    D(q,p;z)\succeq \rho I
\]
on \(\mathcal K\). When \(\rho>0\), the system is strictly dissipative in velocity.
\end{assumption}

\begin{assumption}[Local Lipschitz rollout map]
\label{ass:lipschitz_rollout}
The shared geometry-only update map is locally Lipschitz on \(\mathcal K\): there exists
\(L_\Phi\ge 0\) such that one step of the shared dynamics amplifies state errors by at
most \(1+\tau L_\Phi\).
\end{assumption}

% ======================================================================
% Add these as new subsections within the existing Appendix A
% Slot after existing A.1 (Setup and notation) and before A.3 (Lemma 1)
% ======================================================================

% ----------------------------------------------------------------------
\subsection{Full local-risk interface and problem setup}
\label{app:setup_full}
% ----------------------------------------------------------------------

At each time step $t$ the agent observes a goal displacement
$\vec c_g-\vec c_t$, locally-sensed obstacle or neighbor features
$\{(C_j,R_j,W_j)\}_{j=1}^{N(t)}$ within sensing radius $\hat d$, and a
$P\times P$ context patch $P_t$ centered at the current position.
The patch contains a smoothed risk field $\tilde r(\cdot)\in[0,1]$ and a
signed-distance field $\phi(\cdot)$ to hard hazards or infeasible regions.
Depending on the domain, $\tilde r$ may encode material risk, traffic
interaction risk, or a procedurally constructed local interaction field;
the method only requires that the field be locally queryable and
differentiable after smoothing.
The patch is bilinearly resampled at the agent's current position at
every integration step, so forces and costs reflect the actual path
traversed.
All comparisons use the same local-information regime; differences come
from the decision mechanism, not from privileged sensing.

Each observation defines a local optimization instance.
The goal term, clearance barriers, context-risk preference, hazard
barrier, and lateral-motion channel form a product space of potential
actions and penalties, but only some factors are active in a given
context.
We use \emph{action-space adaptation} in this restricted continuous
sense: the learner does not enumerate discrete actions, but changes which
force channels have non-negligible magnitude under the sensed context.

\paragraph{Static and dynamic risk.}
The same interface covers both static and dynamic settings.
In DFC and RELLIS, $\tilde r$ and $\phi$ are fixed semantic/material
fields, so risk exposure depends on the path taken through the map.
In highway-env, the local patch changes as neighboring vehicles move,
close gaps, or block lateral escape.
Static domains test context-enriched path selection; dynamic domains test
whether the learned field preserves clearance, time-to-collision, and
maneuvering room before violations become terminal.

% ----------------------------------------------------------------------
\subsection{Route-aware gate: full specification}
\label{app:gate}
% ----------------------------------------------------------------------

The main text uses the gate $m_{\rm feas}\in[0,1]$ in
Eq.~\eqref{eq:method:route_aware_force} without specifying its computation.
This subsection gives the full specification.

From the current state $q_t$ we roll out $K$ short-horizon motion
primitives $\{\xi_k\}_{k=1}^K$ inside the sensed BEV window $P_t$ and
compute for each candidate:
\begin{equation}
R_k = \int_{\xi_k}\tilde r(s)\,ds,
\qquad
C_k = \min_{\xi_k}\phi(s),
\qquad
I_k = \mathbf{1}\bigl[C_k>\delta_\phi
      \;\text{and}\;\xi_k\subset\mathcal{T}(P_t)\bigr],
\label{eq:gate_primitives}
\end{equation}
where $\mathcal{T}(P_t)$ is the traversable mask induced by the local
semantics.
Let $k^\star$ be the feasible candidate ($I_{k^\star}=1$) with lowest
soft-risk exposure, and let $k_{\rm scaf}$ be the candidate closest to
the geometry-only policy direction.
The affordance gate is
\begin{equation}
m_{\rm feas}(\vec c,z_t)
=
\sigma\!\bigl(\kappa_R[(R_{k_{\rm scaf}}-R_{k^\star})-\rho_R]\bigr)
\cdot
\sigma\!\bigl(\kappa_C(C_{k^\star}-\delta_\phi)\bigr)
\cdot
I_{k^\star}.
\label{eq:gate_full}
\end{equation}
Thus $m_{\rm feas}$ activates when: (i) a feasible candidate exists
($I_{k^\star}=1$), (ii) that candidate has clearance above threshold
($C_{k^\star}>\delta_\phi$), and (iii) its soft-risk integral is lower
than the geometry-only primitive by margin $\rho_R$.
When any condition fails, $m_{\rm feas}\approx0$ and $F_{\rm soft}$ is
suppressed without affecting $F_{\rm hard}$.

During one Hamiltonian integration substep $m_{\rm feas}$ is treated as
a locally measured context coefficient, exactly like the coefficients
predicted from the risk patch.
Equivalently, the substep uses local energy
$H_{\rm ctx}^{\rm ra}=m_{\rm feas}\lambda_s\tilde r+\lambda_h b_{\rm sp}(\phi)$,
with the gate coefficient frozen until the next BEV update.

\textbf{Implementation notes.}
The route-aware context-enriched field is not a full-map planner or a global-route oracle.
It uses only the local BEV patch, the local SDF, and finite-horizon
primitives within the same sensing radius available to all local-sensing
baselines.
The R1/R2/R3 regime labels are computed externally after the fact to
score whether activation was correct; they are not provided to the policy
at any time.
We use $K{=}8$, $\delta_\phi{=}0.15$\,m, $\rho_R{=}0.05$,
$\kappa_R{=}10$, $\kappa_C{=}20$ in all experiments.

% ----------------------------------------------------------------------
\subsection{Coefficient prediction details}
\label{app:coeff_pred}
% ----------------------------------------------------------------------

The context coefficients $(\lambda_s,\lambda_h)$ are predicted by
sigmoid-bounded heads from a CNN encoding of the risk patch $P_t$ fused
with the existing goal-context token:
\begin{equation}
z_r \leftarrow \psi_\omega(P_t),
\qquad
\lambda_s, \lambda_h \leftarrow \lambda^{\max}\,\sigma\!\bigl(\mathrm{MLP}(z_r,z_g)\bigr),
\end{equation}
where $z_g$ is the goal-context token from the geometry-only policy and $\lambda^{\max}$
is a fixed upper bound on each coefficient.
The geometric heads $(\alpha,\beta,\gamma)$ from the geometry-only policy are frozen,
making $(\lambda_s,\lambda_h)$ the only route by which context can alter
the decision field.
This isolation is what makes the constant-coefficient ablation
(Table~\ref{tab:rellis_selectivity}, constant-coeff row) interpretable:
setting $\lambda_s,\lambda_h$ to fixed scalars removes context-sensitivity
while keeping the force-channel structure, isolating the contribution of
adaptive coefficient prediction.

In highway-env we additionally expose a contextual lateral response
coefficient and a TTC-conditioned interaction channel; in boxed traffic
the TTC channel suppresses lateral activation while in the open
slow-leader case it allows activation.
A $\lambda$-entropy regularizer $L_\lambda$ prevents the sigmoid heads
from saturating.

% ----------------------------------------------------------------------
\subsection{Balanced objective and evaluation metric mapping}
\label{app:jbal}
% ----------------------------------------------------------------------

The full training objective combines the CVaR tail-risk term with a
balanced failure-profile cost:
\begin{equation}
\mathcal{J}_{\rm bal}(\pi)
=
\lambda_{\rm goal}C_{\rm goal}
+\lambda_{\rm risk}C_{\rm risk}
+\lambda_{\rm barrier}C_{\rm barrier}
+\lambda_{\rm control}C_{\rm control}
+\lambda_{\rm smooth}C_{\rm smooth}
+\lambda_{\rm dyn}C_{\rm dyn},
\label{eq:jbal}
\end{equation}
where the dynamic-risk term $C_{\rm dyn}$ is active for interaction
domains (highway-env) and zero otherwise.
The terms are not optimized independently: low risk by stopping, low
control by failing to evade, and short paths through hazards are all
failure modes.
The context-enriched field learns a context-conditioned compromise among these terms.

The one-to-one mapping from $\mathcal{J}_{\rm bal}$ terms to evaluation
metrics is deliberate: goal/progress $\to$ success, progress, path
length; soft risk $\to$ cumulative risk, soft CVaR; barriers $\to$
hard-hazard length, off-road events; control/smoothness $\to$ oscillation,
curvature, jerk; dynamic risk $\to$ clearance, TTC violation,
intervention window.
The metrics are not post-hoc decorations; they measure the same terms the
learned field is balancing.

% ----------------------------------------------------------------------
\subsection{Why expected cost is insufficient}
\label{app:cvar_motivation}
% ----------------------------------------------------------------------

Catastrophic contact, collision, or route failure is a tail event: the
bulk of rollouts may look acceptable while the failures that determine
deployment safety are rare.
Expected-cost training is then dominated by typical-case path quality
and allocates too little gradient to the rollouts where the new context
force channels must become active.

This is not merely a risk-aversion preference.
When training under expected cost, the gradient
$\partial\mathcal{L}/\partial\lambda_s$ receives contributions from all
rollouts, including the majority where $\nabla\tilde r$ is small or
where no risk-reduction opportunity exists.
Gradient signal from safe typical-case rollouts dilutes the signal from
the few rollouts where $-\tau\nabla\tilde r$ should drive a large
$\lambda_s$ update.
CVaR fixes this by restricting gradient to the worst $(1-\alpha)$
fraction: on those rollouts, $-\tau\nabla\tilde r$ contributes non-zero
$\partial q_{t+1}/\partial\lambda_s$ at every step, driving $\lambda_s$
toward values that bias the field away from high-risk regions exactly
where the safety signal lives.

The experimental evidence for this is in Table~\ref{tab:delayed_escape}:
context-enriched expected cost vs.\ route-aware CVaR.
Same architecture, same gate, different objective.
CVaR halves false pre-activation ($0.370\to0.180$), raises suppression
rate ($0.630\to0.820$), and improves navigation success by $30\%$
($0.620\to0.810$).
The R1-only subset (Table~\ref{tab:app_r1_only}) confirms the improvement
holds even when an escape genuinely exists, ruling out the explanation
that CVaR simply avoids acting.

\subsection{Context-enriched update loops and training algorithm}
\label{app:ctx_loops_hierarchy}

The context-enriched field's adaptation hierarchy comprises three concurrent loops that 
operate at different timescales, read different inputs, and update 
different objects.

\textbf{Loop 1 -- segment (within episode, deployment-time).}
A rank-1 secant Gauss--Newton corrector adjusts $(\beta,\gamma)$ and 
the top-$K$ nearest $\alpha_j$ between consecutive frames using realized 
clearance, goal distance, and speed. The Jacobian is maintained as 
$\hat J \leftarrow \mu\hat J + (1{-}\mu)\Delta y\Delta\zeta^\top/
\|\Delta\zeta\|^2$ and a damped least-squares step 
$\Delta\zeta^* = (\hat J^\top\hat J + \rho I)^{-1}\hat J^\top
(y_{\rm tgt}-y)$ is applied with per-coordinate rates and projection to 
$\mathbb{R}_{\ge0}$. This loop runs at every simulator step and uses 
no extra rollouts.

\textbf{Loop 2 -- episode (training-time, between rollouts).}
Backpropagation through the differentiable rollout integrator against 
Eq.~\eqref{eq:method:cvar_emp}, plus auxiliary losses: imitation to a risk-aware 
target ($L_{\rm traj}, L_{\rm vel}$), multi-start robustness 
($L_{\rm multi}$), clearance penalty ($L_{\rm clear}$), 
and $\lambda$-entropy regularizer ($L_\lambda$). Updates the full 
parameter vector including $\theta_R$ and risk encoder $\omega$.

\textbf{Loop 3 -- curriculum (training-time, between phases).}
Geometry-backbone weights are loaded from a geometry-only checkpoint with 
$\lambda_s, \lambda_h \equiv 0$ so $F_{\rm ctx} \equiv 0$ initially. 
Imitation weights are reduced to $0.3\times$ their geometry-only values. 
A passivity surrogate $\mathcal{L}_{\rm pass}$ on a held-out window 
gates phase advancement; the curriculum gradually exposes harder 
material patches, narrower corridors, stronger hazard penalties, and 
interaction scenarios.

\begin{figure}[h]
\centering
\includegraphics[width=\linewidth]{NeurIPS_2026/figures/ch4_three_loops.png}
\caption{\textbf{Three loops as three distinct computational jobs.}
The segment loop corrects active coefficients per step. The episode 
loop optimizes meta-parameters via CVaR. The curriculum loop 
advances training phases statistically.}
\label{fig:app_three_loops}
\end{figure}

\subsection{Context-enriched training algorithm}
\label{app:ctx_training_algorithm}

\begin{algorithm}[t]
\caption{Context-aware continual Hamiltonian learning (training; Loops 2 \& 3)}
\label{alg:ctx_train}
\begin{algorithmic}[1]
\Require geometry-only checkpoint $\theta_{\rm geo} = (\alpha,\beta,\gamma,\theta_R)$;
         dataset $\mathcal{D}$ of short rollouts with local context
         patches; CVaR level $\alpha_c{=}0.95$; batch size $B$;
         horizon $H$; step $\tau$.
\Statex \textbf{Loop 3 -- backward-compatible geometry-only $\to$ context-enriched initialization.}
\State $\theta_{\rm enr} \leftarrow \theta_{\rm scaf}$; initialize risk encoder
        $\omega$ and heads so $\lambda_s, \lambda_h \equiv 0$ at output
       \Comment{$F_{\mathrm{ctx}} \equiv 0$ initially}
\State Reduce imitation weights $(w_{\mathrm{traj}}, w_{\mathrm{vel}})
        \leftarrow 0.3\,(w_{\mathrm{traj}}, w_{\mathrm{vel}})$
\For{epoch $e = 1, \dots, E$}
  \Statex \hspace{0.7em}\textbf{Loop 2 -- episode loop (one minibatch).}
  \For{minibatch $\{(\Pos_0^{(i)}, \Mom_0^{(i)}, \Pos_g^{(i)},
                     \mathcal{O}^{(i)}, P^{(i)}, H^{(i)})\}_{i=1}^{B}
        \in \mathcal{D}$}
    \State $(\alpha^{(i)},\beta^{(i)},\gamma^{(i)})
            \leftarrow \phi_{\mathrm{geom}}(\mathcal{O}^{(i)},
                       \Pos_g^{(i)}{-}\Pos_0^{(i)})$
            \Comment{geometry transformer, unchanged from geometry-only policy}
    \State $z_r^{(i)} \leftarrow \psi_\omega(P^{(i)})$;\;
           $\lambda_s^{(i)},\lambda_h^{(i)} \leftarrow$
           $\lambda^{\max}\sigma(\text{MLP}(z_r^{(i)}, z_g^{(i)}))$
    \For{$t = 0, \dots, H^{(i)}{-}1$ (differentiable rollout)}
      \State $(\tilde r_t,\phi_t,\nabla\tilde r_t,\nabla\phi_t)
              \leftarrow \mathrm{BilinearSample}(P^{(i)},\Pos_t^{(i)})$
              \Comment{resample at \emph{current} $\Pos$, not $\Pos_0$}
      \State $m_{\mathrm{feas}}\leftarrow\mathrm{LocalGate}(P^{(i)},\Pos_t^{(i)},\Pos_g^{(i)})$
              \Comment{Eq.~\eqref{eq:method:route_gate}; set to $1$ for scalar context-enriched field}
      \State $F_{\mathrm{soft}} \leftarrow -m_{\mathrm{feas}}\lambda_s^{(i)} \nabla\tilde r_t$;\;
             $F_{\mathrm{hard}} \leftarrow -\lambda_h^{(i)}\,
                                          b_{\mathrm{sp}}'(\phi_t) \nabla\phi_t$
              \Comment{Eq.~\eqref{eq:method:route_aware_force}}
      \State $\Mom_{t+1}^{(i)} \leftarrow \Mom_t^{(i)}
              + \tau(F_{\mathrm{geom}} + F_{\mathrm{soft}}
                     + F_{\mathrm{hard}})
              - \tau\gamma^{(i)} \Mass^{-1}\Mom_t^{(i)}$
              \Comment{$R_\theta = \gamma\Mass^{-1}\Mom$;
                       $\Mass{=}I$ in implementation}
      \State $\Pos_{t+1}^{(i)} \leftarrow \Pos_t^{(i)} + \tau \Mass^{-1}\Mom_{t+1}^{(i)}$
      \State Accumulate $\mathrm{cum\_risk}, \mathrm{hard\_count}, \mathrm{arc}$
    \EndFor
    \State $J^{(i)} \leftarrow$ Eq.~\eqref{eq:method:J}
    \State $\hat\eta \leftarrow \widehat{Q}_{\alpha_c}(\{J^{(i)}\}_{i=1}^B)$
            \Comment{empirical quantile, \textbf{detach from graph}}
    \State $\widehat{\cvar{\alpha_c}} \leftarrow
           \hat\eta + \tfrac{1}{(1-\alpha_c)B}\sum_i (J^{(i)} - \hat\eta)_+$
            \Comment{Eq.~\eqref{eq:method:cvar_emp}; gradient flows
                     only through tail rollouts}
    \State $\mathcal{L} \leftarrow \widehat{\cvar{\alpha_c}}
                + \sum_{k} w_k L_k$
            \Comment{$L_k \in \{L_{\mathrm{traj}}, L_{\mathrm{vel}},
                                L_{\mathrm{fric}}, L_{\mathrm{clear}},
                                L_{\mathrm{multi}}, L_\lambda\}$}
    \State $\theta_{\rm enr} \leftarrow \theta_{\rm enr} - \mathrm{lr}\,
            \nabla_{\theta_{\rm enr}}\mathcal{L}$,
            $\;\|\nabla\|_2$ clipped to $5.0$
  \EndFor
  \Statex \hspace{0.7em}\textbf{Loop 3 -- curriculum advancement check.}
  \State Compute passivity surrogate $\mathcal{L}_{\mathrm{pass}}$ on
         validation window
  \If{violation rate $< \eta_{\mathrm{tol}}$ over last $W$ epochs}
    \State Mark phase as advanced; optionally widen risk-class
           distribution in $\mathcal{D}$
  \EndIf
\EndFor
\State \Return $\theta_{\rm enr}$
\end{algorithmic}
\end{algorithm}

\subsection{Energy enrichment induces force and sensitivity channels}

\begin{lemma}[Energy enrichment induces force, sensitivity, and excitation channels]
\label{lem:enrichment_channels}
Let
\[
    H_{\rm enr}(q,p;\theta,\vartheta)
    =
    H_{\rm scaf}(q,p;\theta)+\vartheta H_{\mathrm{new}}(q),
\]
where \(H_{\mathrm{new}}\in C^1\) and \(\vartheta\in\mathbb R\) is a learnable scalar.
Then:
\begin{enumerate}
    \item the enriched momentum equation contains the new force channel
    \[
        F_{\mathrm{new}}(q)=-\vartheta\nabla_q H_{\mathrm{new}}(q);
    \]
    \item differentiating a rollout-dependent loss \(J\) through the discrete update
    creates the direct sensitivity stream
    \[
        \frac{\partial J}{\partial \vartheta}
        =
        -\tau\sum_t
        \left\langle
        \frac{\partial J}{\partial p_{t+1}},
        \nabla_qH_{\mathrm{new}}(q_t)
        \right\rangle
        +
        \textnormal{recursive terms};
    \]
    \item local identifiability of \(\vartheta\) from force observations requires
    non-degenerate excitation of the new channel, i.e. the empirical Gram matrix
    over \(\{\nabla_qH_{\mathrm{new}}(q_t)\}_t\) must be nonsingular on the
    observed force subspace.
\end{enumerate}
\end{lemma}

\begin{proof}
Since
\[
    \nabla_q H_{\rm enr}
    =
    \nabla_q H_{\rm scaf}+\vartheta\nabla_qH_{\mathrm{new}},
\]
Hamilton's momentum equation contains
\[
    -\nabla_qH_{\rm enr}
    =
    -\nabla_qH_{\rm scaf}-\vartheta\nabla_qH_{\mathrm{new}}.
\]
Thus the added term contributes the force channel
\(F_{\mathrm{new}}=-\vartheta\nabla_qH_{\mathrm{new}}\).

For the sensitivity statement, consider the discrete momentum update
\[
    p_{t+1}
    =
    p_t-\tau\nabla_qH_{\rm scaf}(q_t,p_t)
    -\tau\vartheta\nabla_qH_{\mathrm{new}}(q_t)+\cdots .
\]
Differentiating with respect to \(\vartheta\) gives
\[
    \frac{\partial p_{t+1}}{\partial\vartheta}
    =
    \frac{\partial p_t}{\partial\vartheta}
    -\tau\nabla_qH_{\mathrm{new}}(q_t)
    -\tau\nabla_q^2H_{\rm enr}(q_t,p_t)
        \frac{\partial q_t}{\partial\vartheta}
    +\cdots .
\]
Applying the chain rule to any rollout loss \(J\) yields
\[
    \frac{\partial J}{\partial\vartheta}
    =
    \sum_t
    \left\langle
    \frac{\partial J}{\partial p_{t+1}},
    \frac{\partial p_{t+1}}{\partial\vartheta}
    \right\rangle
    +
    \sum_t
    \left\langle
    \frac{\partial J}{\partial q_{t+1}},
    \frac{\partial q_{t+1}}{\partial\vartheta}
    \right\rangle .
\]
The displayed direct stream is the term produced by
\(-\tau\nabla_qH_{\mathrm{new}}(q_t)\); the remaining terms are recursive
contributions through \(\partial q_t/\partial\vartheta\) and
\(\partial p_t/\partial\vartheta\).

For identifiability, suppose two coefficients \(\vartheta,\vartheta'\) induce the same
observed force on the rollout. Then
\[
    (\vartheta-\vartheta')\nabla_qH_{\mathrm{new}}(q_t)=0
\]
in all observed force directions. A nonzero coefficient difference is impossible only
when the observed gradients span the relevant force subspace. Equivalently, the
corresponding empirical Gram matrix is nonsingular on that subspace.
\end{proof}

\subsection{Force-dictionary separation}
\label{app:force_dictionary_separation}

\begin{lemma}[Context forces cannot be obtained by geometric reweighting]
Let the geometry-only force dictionary at \(q\) be
\[
    \mathcal D_{\mathrm{geom}}(q)
    =
    \operatorname{span}\{\nabla_q D_{\mathcal G}(q),
    \nabla_q b_{\mathrm{IPC}}(d_1(q)),\ldots,
    \nabla_q b_{\mathrm{IPC}}(d_N(q))\}.
\]
If \(P_{\mathcal D^\perp}\nabla_q\widetilde r(q)\neq0\), then there is no
choice of geometry-only coefficient reweighting that reproduces
\(F_{\mathrm{soft}}=-\lambda_s\nabla_q\widetilde r(q)\) for
\(\lambda_s\neq0\).
\end{lemma}

\begin{proof}
Any geometry-only coefficient reweighting remains in
\(\mathcal D_{\mathrm{geom}}(q)\). But
\(F_{\mathrm{soft}}\) has nonzero projection onto
\(\mathcal D_{\mathrm{geom}}(q)^\perp\) whenever
\(\lambda_sP_{\mathcal D^\perp}\nabla_q\widetilde r(q)\neq0\), so it cannot equal
any vector in the geometry-only span.
\end{proof}

\subsection{Passivity and practical discrete dissipativity}

\begin{theorem}[Contextual Hamiltonian enrichment preserves dissipativity]
\label{thm:continuous_passivity}
Under Assumptions~\ref{ass:regularity}--\ref{ass:diss}, if the context \(z\) is
fixed and \(u=0\), then along solutions of \eqref{eq:ct_q}--\eqref{eq:ct_p},
\[
    \frac{d}{dt}H_{\rm enr}(q(t),p(t);z)
    =
    -v(t)^\top D(q(t),p(t);z)v(t)
    \le
    -\rho\|v(t)\|^2
    \le 0 .
\]
Thus the addition of \(H_{\mathrm{ctx}}\) preserves the port-Hamiltonian dissipativity
structure. If \(z=z(t)\) and \(u\neq0\), then
\[
    \frac{d}{dt}H_{\rm enr}
    =
    -v^\top Dv
    +
    \left\langle \nabla_zH_{\mathrm{ctx}}(q;z),\dot z \right\rangle
    +
    \left\langle \nabla_pH_{\rm enr},Bu \right\rangle .
\]
Hence the enriched system is passive with storage \(H_{\rm enr}\) and supply rate
\[
    s(t)
    =
    \left\langle \nabla_zH_{\mathrm{ctx}}(q;z),\dot z \right\rangle
    +
    \left\langle \nabla_pH_{\rm enr},Bu \right\rangle .
\]
\end{theorem}

\begin{proof}
For fixed \(z\), differentiate \(H_{\rm enr}\) along trajectories:
\[
    \dot H_{\rm enr}
    =
    \langle\nabla_qH_{\rm enr},\dot q\rangle
    +
    \langle\nabla_pH_{\rm enr},\dot p\rangle .
\]
Using \(\dot q=\nabla_pH_{\rm enr}\) and
\(\dot p=-\nabla_qH_{\rm enr}-D\nabla_pH_{\rm enr}\), we obtain
\[
    \dot H_{\rm enr}
    =
    \langle\nabla_qH_{\rm enr},\nabla_pH_{\rm enr}\rangle
    +
    \left\langle
    \nabla_pH_{\rm enr},
    -\nabla_qH_{\rm enr}-D\nabla_pH_{\rm enr}
    \right\rangle .
\]
The two interconnection terms cancel:
\[
    \langle\nabla_qH_{\rm enr},\nabla_pH_{\rm enr}\rangle
    -
    \langle\nabla_pH_{\rm enr},\nabla_qH_{\rm enr}\rangle=0.
\]
Therefore
\[
    \dot H_{\rm enr}
    =
    -\langle\nabla_pH_{\rm enr},D\nabla_pH_{\rm enr}\rangle.
\]
Since \(\nabla_pH_{\rm enr}=M^{-1}p=v\),
\[
    \dot H_{\rm enr}=-v^\top Dv.
\]
By Assumption~\ref{ass:diss}, \(D\succeq \rho I\), so
\[
    -v^\top Dv\le -\rho\|v\|^2\le 0.
\]
If \(z=z(t)\), the chain rule adds
\[
    \langle\nabla_zH_{\rm enr},\dot z\rangle
    =
    \langle\nabla_zH_{\mathrm{ctx}},\dot z\rangle .
\]
If \(u\neq 0\), the input term \(Bu\) in \(\dot p\) contributes
\[
    \langle\nabla_pH_{\rm enr},Bu\rangle .
\]
Combining these terms gives the claimed supply-rate identity.
\end{proof}

\begin{theorem}[Practical discrete dissipativity and step-size condition]
\label{thm:discrete_passivity}
Under Assumptions~\ref{ass:regularity}--\ref{ass:diss}, let \(u_t=0\).
For the semi-implicit context-enriched update \eqref{eq:disc_p}--\eqref{eq:disc_q}, there
exists a finite constant \(C_{\mathcal K}>0\) such that
\[
    H_{\rm enr}(x_{t+1};z_t)-H_{\rm enr}(x_t;z_t)
    \le
    -\tau v_t^\top D_tv_t+C_{\mathcal K}\tau^2 .
\]
Consequently, the discrete rollout is practically dissipative. If
\[
    \tau < \frac{v_t^\top D_tv_t}{C_{\mathcal K}},
\]
then \(H_{\rm enr}(x_{t+1};z_t)<H_{\rm enr}(x_t;z_t)\). If
\(v_t^\top D_tv_t\ge d_{\min}>0\) uniformly on a region of interest, then any
\[
    \tau < \frac{d_{\min}}{C_{\mathcal K}}
\]
guarantees one-step energy decrease on that region.
\end{theorem}

\begin{proof}
Let \(f_z(x)\) be the continuous-time vector field defined by
\eqref{eq:ct_q}--\eqref{eq:ct_p} with \(u=0\). Since the semi-implicit update is a
first-order consistent discretization of this vector field and all derivatives are bounded
on compact \(\mathcal K\), there exists \(C_I>0\) such that
\[
    x_{t+1}=x_t+\tau f_{z_t}(x_t)+e_t,
    \qquad
    \|e_t\|\le C_I\tau^2 .
\]
Because \(H_{\rm enr}\in C^2(\mathcal K)\), its gradient is Lipschitz. Let \(L_H\) be a
Lipschitz constant for \(\nabla H_{\rm enr}\). The smoothness inequality gives
\[
    H_{\rm enr}(x_{t+1})-H_{\rm enr}(x_t)
    \le
    \langle\nabla H_{\rm enr}(x_t),x_{t+1}-x_t\rangle
    +
    \frac{L_H}{2}\|x_{t+1}-x_t\|^2 .
\]
Substituting \(x_{t+1}-x_t=\tau f_{z_t}(x_t)+e_t\),
\begin{align*}
    H_{\rm enr}(x_{t+1})-H_{\rm enr}(x_t)
    &\le
    \tau\langle\nabla H_{\rm enr}(x_t),f_{z_t}(x_t)\rangle
    +
    \langle\nabla H_{\rm enr}(x_t),e_t\rangle  \\
    &\qquad
    +
    \frac{L_H}{2}\|\tau f_{z_t}(x_t)+e_t\|^2 .
\end{align*}
By Theorem~\ref{thm:continuous_passivity},
\[
    \langle\nabla H_{\rm enr}(x_t),f_{z_t}(x_t)\rangle
    =
    -v_t^\top D_tv_t .
\]
The remaining two terms are bounded by \(C_{\mathcal K}\tau^2\) for some finite
constant \(C_{\mathcal K}\), since \(\mathcal K\) is compact and
\(\|\nabla H_{\rm enr}\|\), \(\|f_z\|\), and \(\|e_t\|/\tau^2\) are bounded there. Hence
\[
    H_{\rm enr}(x_{t+1};z_t)-H_{\rm enr}(x_t;z_t)
    \le
    -\tau v_t^\top D_tv_t+C_{\mathcal K}\tau^2 .
\]
Strict decrease follows whenever
\(C_{\mathcal K}\tau^2<\tau v_t^\top D_tv_t\). Dividing by \(\tau>0\) yields the
step-size condition. The uniform-margin result follows by replacing
\(v_t^\top D_tv_t\) by \(d_{\min}\).
\end{proof}

\subsection{Geometry-only preservation and no-hallucination}

\begin{theorem}[Geometry-only preservation under small context force]
\label{thm:scaffold_preservation}
Let \(x_t^{\rm geo}=(q_t^{\rm geo},p_t^{\rm geo})\) be the geometry-only rollout and
\(x_t^{\rm enr}=(q_t^{\rm enr},p_t^{\rm enr})\) be the context-enriched rollout, initialized from the same
state, and suppose Assumptions~\ref{ass:regularity},
\ref{ass:bounded_ctx}, and \ref{ass:lipschitz_rollout} hold. If along the rollout tube
\[
    \|\nabla_qH_{\mathrm{ctx}}(q;z)\|\le \varepsilon .
\]
Then after \(N\) steps,
\[
    \|x_N^{\rm enr}-x_N^{\rm geo}\|
    \le
    \tau\sqrt{1+\tau^2\|M^{-1}\|^2}\,
    \varepsilon\,
    \frac{(1+\tau L_\Phi)^N-1}{\tau L_\Phi},
\]
with the last factor interpreted as \(N\) when \(L_\Phi=0\). In particular, if
\(\varepsilon=0\), the context-enriched rollout recovers the geometry-only rollout up to numerical
roundoff.
\end{theorem}

\begin{proof}
The only difference between the geometry-only and context-enriched fields is the context term in the
momentum update. At one step, this perturbation is
\[
    \delta p_{t+1}
    =
    -\tau\nabla_qH_{\mathrm{ctx}}(q_t;z_t),
\]
and the induced position perturbation is
\[
    \delta q_{t+1}
    =
    \tau M^{-1}\delta p_{t+1}
    =
    -\tau^2M^{-1}\nabla_qH_{\mathrm{ctx}}(q_t;z_t).
\]
Thus
\[
    \|\delta p_{t+1}\|\le \tau\varepsilon,
    \qquad
    \|\delta q_{t+1}\|\le \tau^2\|M^{-1}\|\varepsilon.
\]
The one-step state perturbation is therefore bounded by
\[
    \|\delta x_{t+1}\|
    \le
    \tau\sqrt{1+\tau^2\|M^{-1}\|^2}\,\varepsilon .
\]
Let
\[
    e_t=\|x_t^{\rm enr}-x_t^{\rm geo}\|.
\]
By Assumption~\ref{ass:lipschitz_rollout},
\[
    e_{t+1}
    \le
    (1+\tau L_\Phi)e_t
    +
    \tau\sqrt{1+\tau^2\|M^{-1}\|^2}\,\varepsilon .
\]
Since \(e_0=0\), iterating gives
\[
    e_N
    \le
    \tau\sqrt{1+\tau^2\|M^{-1}\|^2}\,\varepsilon
    \sum_{j=0}^{N-1}(1+\tau L_\Phi)^j .
\]
Evaluating the geometric sum yields the result.
\end{proof}

\begin{corollary}[No hallucinated lateral escape]
\label{cor:no_hallucination}
Let \(P_\perp\) denote projection onto the lateral subspace. Suppose local
coordinates are baseline-aligned so that longitudinal context forces do not generate
lateral motion to first order. If
\[
    \|P_\perp\nabla_qH_{\mathrm{ctx}}(q;z)\|\le \varepsilon_\perp
\]
on the rollout tube, then
\[
    \|P_\perp(q_N^{\rm enr}-q_N^{\rm geo})\|
    \le
    C_{\perp,N}\varepsilon_\perp,
\]
where
\[
    C_{\perp,N}
    =
    \tau^2
    \|P_\perp M^{-1}P_\perp^\top\|
    \frac{(1+\tau L_\Phi)^N-1}{\tau L_\Phi}.
\]
When \(\varepsilon_\perp=0\), the context channel cannot produce a lateral escape.
\end{corollary}

\begin{proof}
Project the context-induced position perturbation:
\[
    P_\perp \delta q_{t+1}
    =
    -\tau^2P_\perp M^{-1}\nabla_qH_{\mathrm{ctx}}(q_t;z_t).
\]
Under baseline-aligned local coordinates, the first-order longitudinal--lateral cross
terms vanish, so
\[
    \|P_\perp \delta q_{t+1}\|
    \le
    \tau^2
    \|P_\perp M^{-1}P_\perp^\top\|
    \varepsilon_\perp .
\]
Propagating this one-step projected perturbation through the same Lipschitz recursion
as in Theorem~\ref{thm:scaffold_preservation} gives the claimed bound.
\end{proof}

\subsection{Selective risk deflection and risk reduction}

\begin{theorem}[Selective risk deflection and local risk reduction]
\label{thm:risk_deflection}
Let \(P_\perp\) project onto a feasible lateral subspace. Suppose that at a state \(q\),
\[
    \|P_\perp\nabla_q\widetilde r(\mathbf c;z)\|\ge \Delta,
\]
and suppose the projected hard-hazard force does not cancel the soft-risk direction by
more than \(\chi\):
\[
    \left\|
    P_\perp
    \left[
    \lambda_h b_{\mathrm{sp}}'(\phi(\mathbf c;z))\nabla_q\phi(\mathbf c;z)
    \right]
    \right\|
    \le \chi.
\]
Then
\[
    \|P_\perp F_{\mathrm{ctx}}(q;z)\|
    \ge
    \lambda_s\Delta-\chi .
\]
Thus, whenever \(\lambda_s\Delta>\chi\), the context channel is necessarily active in
the lateral subspace.

Assume further that \(\widetilde r\) has \(L_r\)-Lipschitz gradient,
\(M^{-1}\succeq m_M I\), and
\[
    \tau^2\lambda_s\|M^{-1}\|\le \frac{1}{L_r}.
\]
Then, ignoring shared geometry and dissipation terms up to \(O(\tau^3)\), the
one-step context-enriched perturbation satisfies
\[
    \widetilde r(q_{t+1}^{\rm enr})
    -
    \widetilde r(q_{t+1}^{\rm scaf})
    \le
    -\frac12\tau^2\lambda_s m_M \Delta^2
    +
    O(\tau^3).
\]
Over any active window \(\mathcal A\) on which the same margin condition holds,
\[
    \sum_{t\in\mathcal A}
    \left[
    \widetilde r(q_{t+1}^{\rm enr})
    -
    \widetilde r(q_{t+1}^{\rm scaf})
    \right]
    \le
    -\frac12|\mathcal A|\tau^2\lambda_s m_M \Delta^2
    +
    O(|\mathcal A|\tau^3).
\]
\end{theorem}

\begin{proof}
The context force is
\[
    F_{\mathrm{ctx}}
    =
    -\lambda_s\nabla_q\widetilde r
    -
    \lambda_h b_{\mathrm{sp}}'(\phi)\nabla_q\phi .
\]
Projecting onto the lateral subspace gives
\[
    P_\perp F_{\mathrm{ctx}}
    =
    -\lambda_sP_\perp\nabla_q\widetilde r
    -
    P_\perp\!\left[
    \lambda_h b_{\mathrm{sp}}'(\phi)\nabla_q\phi
    \right].
\]
By the reverse triangle inequality,
\[
    \|P_\perp F_{\mathrm{ctx}}\|
    \ge
    \lambda_s\|P_\perp\nabla_q\widetilde r\|
    -
    \left\|
    P_\perp\!\left[
    \lambda_h b_{\mathrm{sp}}'(\phi)\nabla_q\phi
    \right]
    \right\|.
\]
The assumptions imply
\[
    \|P_\perp F_{\mathrm{ctx}}\|\ge \lambda_s\Delta-\chi.
\]
This proves activation under \(\lambda_s\Delta>\chi\).

For the risk-reduction claim, isolate the soft-risk force
\[
    F_{\mathrm{soft}}=-\lambda_s\nabla_q\widetilde r(q).
\]
The position perturbation induced by this force in one semi-implicit step is
\[
    \delta q
    =
    -\tau^2\lambda_sM^{-1}\nabla_q\widetilde r(q).
\]
By \(L_r\)-smoothness,
\[
    \widetilde r(q+\delta q)
    \le
    \widetilde r(q)
    +
    \langle\nabla_q\widetilde r(q),\delta q\rangle
    +
    \frac{L_r}{2}\|\delta q\|^2 .
\]
The first-order term is
\[
    \langle\nabla_q\widetilde r(q),\delta q\rangle
    =
    -\tau^2\lambda_s
    \nabla_q\widetilde r(q)^\top
    M^{-1}
    \nabla_q\widetilde r(q).
\]
Since \(M^{-1}\succeq m_MI\) and
\(\|\nabla_q\widetilde r(q)\|\ge \Delta\) on the active projected direction,
\[
    \nabla_q\widetilde r(q)^\top M^{-1}\nabla_q\widetilde r(q)
    \ge
    m_M\Delta^2.
\]
Therefore
\[
    \langle\nabla_q\widetilde r(q),\delta q\rangle
    \le
    -\tau^2\lambda_s m_M\Delta^2.
\]
The second-order term is
\[
    \frac{L_r}{2}\|\delta q\|^2
    =
    \frac{L_r}{2}
    \tau^4\lambda_s^2
    \|M^{-1}\nabla_q\widetilde r(q)\|^2 .
\]
The step-size condition
\[
    \tau^2\lambda_s\|M^{-1}\|\le \frac{1}{L_r}
\]
ensures that this second-order term is controlled by the first-order decrease, yielding
\[
    \widetilde r(q+\delta q)-\widetilde r(q)
    \le
    -\frac12\tau^2\lambda_s m_M\Delta^2.
\]
The shared geometry and dissipative terms are identical in the geometry-only and context-enriched fields
local comparison; by smoothness, their contribution to the relative risk difference is
\(O(\tau^3)\). Summing the one-step inequality over \(\mathcal A\) proves the final
claim.
\end{proof}

\subsection{CVaR gradient consistency}

The episode cost used in the main text is
\[
    J(\theta)
    =
    w_g\|q_T-q_g\|^2
    +
    w_\ell\sum_t\|q_{t+1}-q_t\|
    +
    w_r\sum_t \widetilde r(q_t)\|q_{t+1}-q_t\|
    +
    w_h\sum_t \mathbf 1\{\phi(q_t)<\epsilon\}.
\]
The training objective uses the Rockafellar--Uryasev representation of CVaR:
\[
    \operatorname{CVaR}_\alpha(J_\theta)
    =
    \min_{\eta\in\mathbb R}
    \left[
    \eta+
    \frac{1}{1-\alpha}\mathbb E(J_\theta-\eta)_+
    \right].
\]

\begin{theorem}[Detached empirical CVaR gives a consistent tail-gradient estimator]
\label{thm:cvar_consistency}
Let \(J_\theta(\xi)\) be the rollout cost induced by parameter \(\theta\) and rollout
randomness \(\xi\). Fix \(\alpha\in(0,1)\). Assume:
\begin{enumerate}
    \item \(J_\theta(\xi)\) is differentiable in \(\theta\) almost surely;
    \item \(\|\nabla_\theta J_\theta(\xi)\|\le G\) almost surely;
    \item the distribution of \(J_\theta\) has a continuous density near its
    \(\alpha\)-quantile;
    \item \(\mathbb P[J_\theta=\eta_\alpha(\theta)]=0\), where
    \(\eta_\alpha(\theta)=\operatorname{VaR}_\alpha(J_\theta)\).
\end{enumerate}
Then
\[
    \nabla_\theta \operatorname{CVaR}_\alpha(J_\theta)
    =
    \frac{1}{1-\alpha}
    \mathbb E
    \left[
    \nabla_\theta J_\theta(\xi)\,
    \mathbf 1\{J_\theta(\xi)\ge \eta_\alpha(\theta)\}
    \right].
\]
Given \(B\) i.i.d. rollouts and the detached empirical quantile
\[
    \widehat \eta
    =
    \widehat Q_\alpha(\{J_\theta^{(i)}\}_{i=1}^B),
\]
define
\[
    \widehat g_B(\theta)
    =
    \frac{1}{(1-\alpha)B}
    \sum_{i=1}^B
    \nabla_\theta J_\theta^{(i)}
    \mathbf 1\{J_\theta^{(i)}\ge \widehat \eta\}.
\]
Then
\[
    \widehat g_B(\theta)
    \to
    \nabla_\theta \operatorname{CVaR}_\alpha(J_\theta)
\]
almost surely as \(B\to\infty\), and
\[
    \|\widehat g_B(\theta)
    -
    \nabla_\theta \operatorname{CVaR}_\alpha(J_\theta)\|
    =
    O_p(B^{-1/2}).
\]
If \(J_\theta\in[0,J_{\max}]\), then the empirical CVaR value obeys
\[
    |\widehat{\operatorname{CVaR}}_\alpha-\operatorname{CVaR}_\alpha|
    =
    O\!\left(
    \frac{J_{\max}}{1-\alpha}
    \sqrt{\frac{\log(1/\delta)}{B}}
    \right)
\]
with probability at least \(1-\delta\).
\end{theorem}

\begin{proof}
The Rockafellar--Uryasev variational form gives
\[
    \operatorname{CVaR}_\alpha(J_\theta)
    =
    \min_{\eta\in\mathbb R}
    \left[
    \eta+
    \frac{1}{1-\alpha}\mathbb E(J_\theta-\eta)_+
    \right].
\]
At the minimizer,
\[
    \eta^\star=\eta_\alpha(\theta).
\]
By the envelope theorem, when differentiating the optimized value with respect to
\(\theta\), the derivative of the optimizer \(\eta^\star(\theta)\) does not appear. Hence
\[
    \nabla_\theta \operatorname{CVaR}_\alpha(J_\theta)
    =
    \frac{1}{1-\alpha}
    \nabla_\theta
    \mathbb E[(J_\theta-\eta^\star)_+].
\]
Since \(J_\theta\) is differentiable almost surely and the quantile boundary has
probability zero,
\[
    \nabla_\theta(J_\theta-\eta^\star)_+
    =
    \nabla_\theta J_\theta\,
    \mathbf 1\{J_\theta\ge \eta^\star\}
\]
almost surely. The bounded-gradient assumption permits interchanging derivative and
expectation, proving
\[
    \nabla_\theta \operatorname{CVaR}_\alpha(J_\theta)
    =
    \frac{1}{1-\alpha}
    \mathbb E
    \left[
    \nabla_\theta J_\theta
    \mathbf 1\{J_\theta\ge \eta_\alpha\}
    \right].
\]

For the empirical estimator, continuity of the distribution near \(\eta_\alpha\) implies
strong consistency of the empirical quantile:
\[
    \widehat \eta \to \eta_\alpha
\]
almost surely. Since there is no atom at the quantile,
\[
    \mathbf 1\{J_\theta^{(i)}\ge\widehat\eta\}
    \to
    \mathbf 1\{J_\theta^{(i)}\ge\eta_\alpha\}
\]
almost surely. The summands are bounded by \(G/(1-\alpha)\), so dominated
convergence and the strong law of large numbers give
\[
    \widehat g_B(\theta)
    \to
    \frac{1}{1-\alpha}
    \mathbb E
    \left[
    \nabla_\theta J_\theta
    \mathbf 1\{J_\theta\ge\eta_\alpha\}
    \right]
    =
    \nabla_\theta \operatorname{CVaR}_\alpha(J_\theta)
\]
almost surely.

For the rate, decompose
\[
    \widehat g_B-g
    =
    A_B+R_B,
\]
where
\[
    A_B
    =
    \frac{1}{(1-\alpha)B}
    \sum_{i=1}^B
    \left[
    \nabla J_i\mathbf 1\{J_i\ge\eta_\alpha\}
    -
    \mathbb E(\nabla J\mathbf 1\{J\ge\eta_\alpha\})
    \right]
\]
and
\[
    R_B
    =
    \frac{1}{(1-\alpha)B}
    \sum_{i=1}^B
    \nabla J_i
    \left[
    \mathbf 1\{J_i\ge\widehat\eta\}
    -
    \mathbf 1\{J_i\ge\eta_\alpha\}
    \right].
\]
The first term is a bounded empirical average, so \(A_B=O_p(B^{-1/2})\). The
second term is nonzero only when \(J_i\) lies between \(\widehat\eta\) and
\(\eta_\alpha\). The probability mass of this interval is
\(O_p(|\widehat\eta-\eta_\alpha|)\), and standard empirical quantile theory gives
\(|\widehat\eta-\eta_\alpha|=O_p(B^{-1/2})\). Hence \(R_B=O_p(B^{-1/2})\), proving
the gradient rate.

For bounded \(J_\theta\), the empirical Rockafellar--Uryasev objective is a bounded
empirical process indexed by \(\eta\). Uniform concentration over
\(\eta\in[0,J_{\max}]\), followed by comparison at the empirical and population
minimizers, yields the displayed
\[
    O\!\left(
    \frac{J_{\max}}{1-\alpha}
    \sqrt{\frac{\log(1/\delta)}{B}}
    \right)
\]
concentration rate.
\end{proof}


\subsection{Risk-averse constrained OCP interpretation}
\label{app:risk_constrained_ocp}

This subsection formalizes the statement that the contextual Hamiltonian should be
viewed as an amortized approximation to a \emph{family} of risk-averse constrained
optimal control problems, not as an unconditional online optimizer. The Hamiltonian
policy is a structured feedback parameterization, while the constrained OCP is the
reference problem whose local Lagrangian force field it aims to approximate.

Consider a finite-horizon controlled system
\begin{equation}
    x_{t+1}=\Phi(x_t,u_t,z_t),
    \qquad x_t=(q_t,p_t),
    \qquad t=0,\dots,H-1,
    \label{eq:app:discrete_ocp_dynamics}
\end{equation}
with local context process $z_t$ and rollout randomness $\xi$. Let $J_0(\tau)$ denote
the nominal goal, path-length, smoothness, and control-effort cost. Let
$C_j(\tau)$ denote the exposure accumulated by risk channel $j$, for example
soft-material exposure, hard-hazard proximity, deformation risk, or interaction risk.
The risk-averse constrained reference family is
\begin{equation}
\begin{aligned}
    \min_{\pi\in\Pi}\quad
    & \mathbb E[J_0(\tau_\pi)] \\
    \mathrm{s.t.}\quad
    & \operatorname{CVaR}_{\alpha_j}(C_j(\tau_\pi))
      \le \bar\rho_j,
      \qquad j=1,\dots,m .
\end{aligned}
\label{eq:app:risk_constrained_ocp}
\end{equation}
Equivalently, for nonnegative multipliers $\lambda\in\mathbb R_+^m$, its Lagrangian
risk objective is
\begin{equation}
    \mathcal J_\lambda(\pi)
    =
    \mathbb E[J_0(\tau_\pi)]
    +
    \sum_{j=1}^m
    \lambda_j
    \bigl(
      \operatorname{CVaR}_{\alpha_j}(C_j(\tau_\pi))-\bar\rho_j
    \bigr).
    \label{eq:app:risk_lagrangian_ocp}
\end{equation}
Under a local constraint qualification, a locally optimal constrained policy admits
KKT multipliers~\citep{ShapiroDentchevaRuszczynski2009} $\lambda^\star$; the corresponding first-order conditions are those of
\eqref{eq:app:risk_lagrangian_ocp} at $\lambda^\star$ together with complementary
slackness. This is the precise sense in which the coefficients of risk-aware energy
terms can be interpreted as local risk prices.

The learned Hamiltonian policy class is
\begin{equation}
    H_\theta(q,p;z)
    =
    \frac12p^\top M^{-1}p
    +U_\theta(q,z),
    \qquad
    u_{\rm raw}=\pi_\theta(x,z),
    \label{eq:app:hamiltonian_policy_class}
\end{equation}
where the induced closed-loop vector field has the form
\begin{equation}
    F_\theta(x,z)
    =
    (J-R_\theta)\nabla_xH_\theta(x,z)+G(x)u_{\rm raw}.
\end{equation}
The next theorem states the approximation claim used in the main text. It is a local
inverse-optimality statement: if the Hamiltonian force approximates the local optimal
Lagrangian closed-loop force, then the resulting trajectories and risk-Lagrangian
values are close.

\begin{theorem}[Hamiltonian approximation of the risk-averse Lagrangian OCP]
\label{thm:hamiltonian_approx_risk_ocp}
Fix a compact rollout tube $\mathcal K$ and a multiplier vector
$\lambda\in\mathbb R_+^m$. Suppose that a locally optimal feedback
$\pi_\lambda^\star$ for \eqref{eq:app:risk_lagrangian_ocp} induces a closed-loop update
map $\Phi_\lambda^\star$ on $\mathcal K$, and the Hamiltonian feedback
$\pi_\theta$ induces a closed-loop update map $\Phi_\theta$. Assume:
\begin{enumerate}
    \item $\Phi_\lambda^\star$ is $L_\Phi$-Lipschitz on $\mathcal K$;
    \item the one-step model mismatch is uniformly bounded,
    \begin{equation}
        \sup_{x,z\in\mathcal K}
        \|\Phi_\theta(x,z)-\Phi_\lambda^\star(x,z)\|
        \le \varepsilon_\Phi ;
        \label{eq:app:one_step_phi_mismatch}
    \end{equation}
    \item the rollout costs $J_0,C_1,\dots,C_m$ are $L_J,L_{C_1},\dots,L_{C_m}$
    Lipschitz functions of the finite trajectory under the sup norm.
\end{enumerate}
Then trajectories initialized at the same $x_0$ satisfy
\begin{equation}
    \max_{0\le t\le H}\|x_t^\theta-x_t^\star\|
    \le
    C_H\varepsilon_\Phi,
    \qquad
    C_H:=\sum_{k=0}^{H-1}(1+L_\Phi)^k .
    \label{eq:app:ocp_traj_close}
\end{equation}
Moreover,
\begin{equation}
\begin{aligned}
    \mathcal J_\lambda(\pi_\theta)-\mathcal J_\lambda(\pi_\lambda^\star)
    \le
    \left(L_J+\sum_{j=1}^m\lambda_jL_{C_j}\right)
    C_H\varepsilon_\Phi .
\end{aligned}
\label{eq:app:risk_lagrangian_close}
\end{equation}
If the constrained problem \eqref{eq:app:risk_constrained_ocp} satisfies a local
constraint qualification and $\lambda^\star$ is a KKT multiplier at a local constrained
solution, then \eqref{eq:app:risk_lagrangian_close} applies to the local Lagrangian
relaxation of the constrained risk-averse OCP at $\lambda^\star$.
\end{theorem}

\begin{proof}
Let $e_t=x_t^\theta-x_t^\star$. By adding and subtracting
$\Phi_\lambda^\star(x_t^\theta,z_t)$,
\[
\begin{aligned}
    \|e_{t+1}\|
    &\le
    \|\Phi_\theta(x_t^\theta,z_t)-\Phi_\lambda^\star(x_t^\theta,z_t)\|
    +
    \|\Phi_\lambda^\star(x_t^\theta,z_t)-\Phi_\lambda^\star(x_t^\star,z_t)\| \\
    &\le
    \varepsilon_\Phi+(1+L_\Phi)\|e_t\| .
\end{aligned}
\]
Since $e_0=0$, discrete Gronwall gives
\[
    \max_{0\le t\le H}\|e_t\|
    \le
    \varepsilon_\Phi\sum_{k=0}^{H-1}(1+L_\Phi)^k
    =C_H\varepsilon_\Phi .
\]
The Lipschitz assumptions on the trajectory functionals imply
\[
    |J_0(\tau_\theta)-J_0(\tau_\star)|
    \le L_JC_H\varepsilon_\Phi,
    \qquad
    |C_j(\tau_\theta)-C_j(\tau_\star)|
    \le L_{C_j}C_H\varepsilon_\Phi .
\]
CVaR is monotone and translation equivariant. Therefore, if two random variables
$X,Y$ obey $|X-Y|\le a$ almost surely, then
$|\operatorname{CVaR}_\alpha(X)-\operatorname{CVaR}_\alpha(Y)|\le a$. Applying this
with $X=C_j(\tau_\theta)$, $Y=C_j(\tau_\star)$ and
$a=L_{C_j}C_H\varepsilon_\Phi$ gives
\[
    \left|
    \operatorname{CVaR}_{\alpha_j}(C_j(\tau_\theta))
    -
    \operatorname{CVaR}_{\alpha_j}(C_j(\tau_\star))
    \right|
    \le L_{C_j}C_H\varepsilon_\Phi .
\]
Combining the nominal and risk-channel bounds yields
\eqref{eq:app:risk_lagrangian_close}. Finally, under the stated constraint
qualification, the KKT conditions identify the local constrained solution with a
stationary point of the Lagrangian relaxation at some nonnegative multiplier
$\lambda^\star$ satisfying complementary slackness; substituting
$\lambda=\lambda^\star$ gives the constrained-OCP interpretation.
\end{proof}

\subsection{Identifiability, generalization, and upper-confidence risk safety}
\label{app:identifiability_generalization_safety}

This subsection makes precise what can and cannot be identified from learning the
risk-aware Hamiltonian, and how the empirical CVaR update becomes a genuine
risk-constraint update once an upper-confidence radius is added. The argument uses
three standard ingredients: uniform convergence for bounded function classes via
Rademacher complexity~\citep{bartlett2002rademacher,ShalevShwartzBenDavid2014},
quantile/CVaR concentration via the Dvoretzky--Kiefer--Wolfowitz inequality with
Massart's constant~\citep{DvoretzkyKieferWolfowitz1956,Massart1990}, and
forward-invariance of barrier-filtered dynamics via Nagumo/CBF conditions
\citep{Nagumo1942,ames2016control}.

\paragraph{Structured risk-energy class.}
For a local context patch $z$, write the learned potential as
\begin{equation}
    U_\theta(q,z)
    =
    U_{\rm goal}(q,z)
    +
    \sum_{i=1}^{N(z)} \alpha_i(z) B_i(q,z)
    +
    U_{{\rm res},\theta}(q,z),
    \label{eq:app:structured_energy}
\end{equation}
where $B_i$ are fixed differentiable obstacle or risk-barrier atoms and
$U_{{\rm res},\theta}$ is the learned residual shaping potential. The induced force is
\begin{equation}
    F_\theta(q,z)
    =
    -\nabla_q U_\theta(q,z).
\end{equation}
The scalar potential is never identifiable beyond an additive constant; only its
gradient can be identified from local force or direction observations. Moreover, the
split between the structured barrier part and the residual is identifiable only under an
excitation and non-absorption condition.

\begin{assumption}[Barrier excitation and residual non-absorption]
\label{ass:barrier_excitation}
For each context patch $z$, let
\[
    G_B(q,z)
    =
    [\nabla_q B_1(q,z),\dots,\nabla_q B_{N(z)}(q,z)]\in\mathbb R^{d\times N(z)}.
\]
There exists $\kappa_B>0$ such that
\begin{equation}
    \lambda_{\min}
    \left(
    \mathbb E_{q\sim\mu_z}[G_B(q,z)^\top G_B(q,z)]
    \right)
    \ge \kappa_B .
    \label{eq:app:barrier_gram}
\end{equation}
The residual force is either constrained to be small,
$\mathbb E_{\mu_z}\|\nabla_qU_{{\rm res},\theta}\|^2\le \varepsilon_{\rm res}^2$, or
orthogonal to the barrier dictionary,
\begin{equation}
    \mathbb E_{q\sim\mu_z}
    \left[G_B(q,z)^\top \nabla_qU_{{\rm res},\theta}(q,z)\right]=0 .
\end{equation}
\end{assumption}

\begin{proposition}[Local identifiability of the structured risk weights]
\label{prop:local_identifiability}
Fix $z$ and suppose the teacher force admits the decomposition
\begin{equation}
    F^\star(q,z)
    =
    -\nabla_q U_{\rm goal}(q,z)
    -G_B(q,z)\alpha^\star(z)
    +\epsilon(q,z),
    \label{eq:app:teacher_force_decomp}
\end{equation}
with $\mathbb E\|\epsilon(q,z)\|^2\le \sigma_\epsilon^2$. Let
$\widehat\alpha(z)$ minimize the population least-squares force error over the
structured class with residual either removed or constrained as in
Assumption~\ref{ass:barrier_excitation}. Then
\begin{equation}
    \|\widehat\alpha(z)-\alpha^\star(z)\|
    \le
    \frac{1}{\kappa_B}
    \left(
    \left\|
    \mathbb E[G_B(q,z)^\top\epsilon(q,z)]
    \right\|
    +
    \varepsilon_{\rm res}
    \sqrt{\mathbb E\|G_B(q,z)\|_{\rm op}^2}
    \right).
    \label{eq:app:alpha_identifiability_bound}
\end{equation}
In particular, if the model is well specified, the residual is orthogonal to the
barrier dictionary, and $\mathbb E[G_B^\top\epsilon]=0$, then
$\widehat\alpha(z)=\alpha^\star(z)$. Without the Gram condition
\eqref{eq:app:barrier_gram}, the structured weights are not identifiable.
\end{proposition}

\begin{proof}
Subtract $-\nabla U_{\rm goal}$ from both the teacher force and the model force. The
population normal equation for the structured coefficients is
\[
    \mathbb E[G_B^\top G_B]\widehat\alpha
    =
    \mathbb E[G_B^\top(G_B\alpha^\star-\epsilon-r_\theta)],
\]
where $r_\theta=\nabla_qU_{{\rm res},\theta}$, with the sign absorbed consistently in
both sides. Hence
\[
    \mathbb E[G_B^\top G_B](\widehat\alpha-\alpha^\star)
    =
    -\mathbb E[G_B^\top\epsilon]-\mathbb E[G_B^\top r_\theta].
\]
Taking norms and using
$\|A^{-1}\|_{\rm op}\le 1/\kappa_B$ gives
\[
    \|\widehat\alpha-\alpha^\star\|
    \le
    \frac{1}{\kappa_B}
    \left(
    \|\mathbb E[G_B^\top\epsilon]\|
    +
    \|\mathbb E[G_B^\top r_\theta]\|
    \right).
\]
If the residual is orthogonal to the dictionary, the second term vanishes. If it is only
small, Cauchy--Schwarz gives
\[
    \|\mathbb E[G_B^\top r_\theta]\|
    \le
    \sqrt{\mathbb E\|G_B\|_{\rm op}^2}
    \sqrt{\mathbb E\|r_\theta\|^2}
    \le
    \varepsilon_{\rm res}
    \sqrt{\mathbb E\|G_B\|_{\rm op}^2}.
\]
This proves the bound. If $\mathbb E[G_B^\top G_B]$ is singular, there exists a
nonzero vector $a$ in its nullspace. Then
$G_B(q,z)a=0$ for $\mu_z$-almost every $q$, so $\alpha$ and $\alpha+a$ induce the
same observed structured force. Thus the coefficients are not identifiable.
\end{proof}

\begin{assumption}[Bounded force class]
\label{ass:bounded_force_class}
Let
\[
    \mathcal F=\{(q,p,z)\mapsto F_\theta(q,p,z):\theta\in\Theta\}
\]
be the force class induced by the Hamiltonian network on the compact rollout tube
$\mathcal K$. Assume that
$\|F_\theta(q,p,z)-F^\star(q,p,z)\|\le M_F$ and that the squared loss class
\[
    \mathcal L_F
    =
    \{(q,p,z,F^\star)\mapsto \|F_\theta(q,p,z)-F^\star\|^2:\theta\in\Theta\}
\]
is measurable and has empirical Rademacher complexity
$\mathfrak R_n(\mathcal L_F)$.
\end{assumption}

\begin{theorem}[Force-field generalization]
\label{thm:force_generalization}
Let $(q_i,p_i,z_i,F_i^\star)_{i=1}^n$ be i.i.d. samples from the local patch
distribution, and define
\[
    \mathcal R_F(\theta)
    =
    \mathbb E\|F_\theta(q,p,z)-F^\star(q,p,z)\|^2,
    \qquad
    \widehat{\mathcal R}_{F,n}(\theta)
    =
    \frac1n\sum_{i=1}^n
    \|F_\theta(q_i,p_i,z_i)-F_i^\star\|^2 .
\]
Under Assumption~\ref{ass:bounded_force_class}, with probability at least
$1-\delta$,
\begin{equation}
    \mathcal R_F(\hat\theta)
    \le
    \widehat{\mathcal R}_{F,n}(\hat\theta)
    +
    2\mathfrak R_n(\mathcal L_F)
    +
    3M_F^2\sqrt{\frac{\log(2/\delta)}{2n}}
    \label{eq:app:force_gen_bound}
\end{equation}
for every data-dependent estimator $\hat\theta\in\Theta$. If the learned dynamics and
costs are Lipschitz on $\mathcal K$, this force error propagates to rollout-cost error
with a finite-horizon Gronwall factor.
\end{theorem}

\begin{proof}
The first statement is the standard Rademacher uniform convergence inequality for a
bounded loss class~\citep{bartlett2002rademacher,ShalevShwartzBenDavid2014}:
with probability at least $1-\delta$,
\[
    \sup_{\theta\in\Theta}
    \bigl(\mathcal R_F(\theta)-\widehat{\mathcal R}_{F,n}(\theta)\bigr)
    \le
    2\mathfrak R_n(\mathcal L_F)
    +
    3M_F^2\sqrt{\frac{\log(2/\delta)}{2n}} .
\]
Substituting the random estimator $\hat\theta$ into this uniform event gives
\eqref{eq:app:force_gen_bound}. For the rollout statement, let
$e_t=x_t^\theta-x_t^\star$. If the closed-loop update map is $L_\Phi$-Lipschitz and
the one-step force mismatch is bounded by $\varepsilon_F$, then
\[
    \|e_{t+1}\|
    \le
    (1+\tau L_\Phi)\|e_t\|+\tau\varepsilon_F .
\]
Iterating gives
\[
    \max_{0\le t\le H}\|e_t\|
    \le
    \tau\varepsilon_F
    \sum_{k=0}^{H-1}(1+\tau L_\Phi)^k .
\]
Lipschitz continuity of the finite-horizon cost then converts the trajectory bound into
an objective bound.
\end{proof}

\begin{lemma}[Upper-confidence CVaR bound]
\label{lem:cvar_ucb}
Let $C(\tau_\theta)\in[0,C_{\max}]$ be a bounded rollout risk exposure and let
$C_1,\dots,C_B$ be i.i.d. validation rollouts for a fixed policy $\theta$. Define the
population and empirical Rockafellar--Uryasev objectives
\begin{align}
    \varphi(\eta)
    &=
    \eta+\frac{1}{1-\alpha}\mathbb E[(C-\eta)_+], \\
    \widehat\varphi_B(\eta)
    &=
    \eta+\frac{1}{(1-\alpha)B}\sum_{i=1}^B(C_i-\eta)_+ .
\end{align}
Let
\[
    \rho_\alpha=\inf_{\eta\in\mathbb R}\varphi(\eta),
    \qquad
    \widehat\rho_{\alpha,B}=\inf_{\eta\in\mathbb R}\widehat\varphi_B(\eta).
\]
Then, with probability at least $1-\delta$,
\begin{equation}
    \rho_\alpha
    \le
    \widehat\rho_{\alpha,B}
    +
    \frac{C_{\max}}{1-\alpha}
    \sqrt{\frac{\log(2/\delta)}{2B}} .
    \label{eq:app:cvar_ucb_single}
\end{equation}
For $m$ bounded risk channels $C_j\in[0,C_{j,\max}]$, the bound holds
simultaneously for all channels with probability at least $1-\delta$ after using
\begin{equation}
    \beta_{j,B}(\delta/m)
    :=
    \frac{C_{j,\max}}{1-\alpha_j}
    \sqrt{\frac{\log(2m/\delta)}{2B}} .
    \label{eq:app:explicit_cvar_beta}
\end{equation}
\end{lemma}

\begin{proof}
Because $C\in[0,C_{\max}]$, every minimizer of $\varphi$ and
$\widehat\varphi_B$ may be chosen in $[0,C_{\max}]$: if $\eta<0$, increasing
$\eta$ to $0$ cannot increase the Rockafellar--Uryasev objective, and if
$\eta>C_{\max}$, decreasing $\eta$ to $C_{\max}$ cannot increase it. Hence it is
enough to control the objectives uniformly over $\eta\in[0,C_{\max}]$.

Let $F$ and $F_B$ denote the population and empirical distribution functions of
$C$. For $\eta\in[0,C_{\max}]$, the tail-integral identity gives
\begin{equation}
    \mathbb E[(C-\eta)_+]
    =
    \int_\eta^{C_{\max}}(1-F(t))\,dt,
    \qquad
    \frac1B\sum_{i=1}^B(C_i-\eta)_+
    =
    \int_\eta^{C_{\max}}(1-F_B(t))\,dt .
    \label{eq:app:tail_integral_identity}
\end{equation}
Therefore, for every $\eta\in[0,C_{\max}]$,
\begin{align}
    |\widehat\varphi_B(\eta)-\varphi(\eta)|
    &\le
    \frac{1}{1-\alpha}
    \int_\eta^{C_{\max}}|F_B(t)-F(t)|\,dt \\
    &\le
    \frac{C_{\max}}{1-\alpha}
    \|F_B-F\|_\infty .
    \label{eq:app:ru_uniform_from_dkw}
\end{align}
Massart's sharp form of the Dvoretzky--Kiefer--Wolfowitz inequality states that
\citep{DvoretzkyKieferWolfowitz1956,Massart1990}
\begin{equation}
    \mathbb P\!\left(\|F_B-F\|_\infty>\epsilon\right)
    \le
    2\exp(-2B\epsilon^2).
\end{equation}
Thus, with probability at least $1-\delta$,
\[
    \|F_B-F\|_\infty
    \le
    \sqrt{\frac{\log(2/\delta)}{2B}} .
\]
Combining this event with \eqref{eq:app:ru_uniform_from_dkw} gives
\begin{equation}
    \sup_{\eta\in[0,C_{\max}]}
    |\widehat\varphi_B(\eta)-\varphi(\eta)|
    \le
    \frac{C_{\max}}{1-\alpha}
    \sqrt{\frac{\log(2/\delta)}{2B}} .
    \label{eq:app:ru_uniform_event}
\end{equation}
Finally, for any two functions $f,g$, $|\inf f-\inf g|\le\sup|f-g|$. Applying this
with $f=\varphi$ and $g=\widehat\varphi_B$ on the event
\eqref{eq:app:ru_uniform_event} yields
\[
    \rho_\alpha
    \le
    \widehat\rho_{\alpha,B}
    +
    \frac{C_{\max}}{1-\alpha}
    \sqrt{\frac{\log(2/\delta)}{2B}} .
\]
The simultaneous result follows by applying the same argument to each channel with
failure probability $\delta/m$ and taking a union bound.
\end{proof}

\begin{theorem}[Upper-confidence risk-constrained learning]
\label{thm:ucb_risk_constrained_learning}
For each risk channel $j=1,\dots,m$, let
\[
    \rho_j(\theta)
    =
    \operatorname{CVaR}_{\alpha_j}(C_j(\tau_\theta)),
    \qquad
    \widehat\rho_{j,B}(\theta)
    =
    \widehat{\operatorname{CVaR}}_{\alpha_j,B}(C_j(\tau_\theta))
    +
    \beta_{j,B}(\delta/m),
\]
where $\beta_{j,B}$ is chosen as in Eq.~\eqref{eq:app:explicit_cvar_beta}. Let
$\hat\theta$ be fixed before drawing the validation rollouts used to compute
$\widehat\rho_{j,B}$; alternatively, let $\beta_{j,B}$ include a uniform
function-class complexity term. If
\begin{equation}
    \widehat\rho_{j,B}(\hat\theta)\le \bar\rho_j,
    \qquad j=1,\dots,m,
    \label{eq:app:empirical_ucb_feasible}
\end{equation}
then, with probability at least $1-\delta$,
\begin{equation}
    \rho_j(\hat\theta)
    \le
    \bar\rho_j,
    \qquad j=1,\dots,m .
    \label{eq:app:true_risk_feasible}
\end{equation}
Moreover, the projected dual update
\begin{equation}
    \lambda_j^{k+1}
    =
    \left[
    \lambda_j^k
    +
    \eta_\lambda
    \bigl(\widehat\rho_{j,B}(\theta_k)-\bar\rho_j\bigr)
    \right]_+
    \label{eq:app:dual_update}
\end{equation}
performs ascent on the empirical upper-bound Lagrangian constraint violation and
therefore increases the risk price exactly when the upper-confidence budget is
violated.
\end{theorem}

\begin{proof}
Because the validation rollouts are independent of the fixed iterate $\hat\theta$,
Lemma~\ref{lem:cvar_ucb} applies to that iterate. A union bound over the $m$ risk
channels implies that, with probability at least $1-\delta$,
\[
    \rho_j(\hat\theta)
    \le
    \widehat\rho_{j,B}(\hat\theta)
\]
holds simultaneously for all $j$. Combining this event with
\eqref{eq:app:empirical_ucb_feasible} gives
\[
    \rho_j(\hat\theta)
    \le
    \widehat\rho_{j,B}(\hat\theta)
    \le
    \bar\rho_j,
\]
which proves true risk feasibility. The dual statement follows because the Lagrangian
term for channel $j$ is
$\lambda_j(\widehat\rho_{j,B}(\theta)-\bar\rho_j)$ with the constraint
$\lambda_j\ge 0$; projected gradient ascent in $\lambda_j$ is exactly
\eqref{eq:app:dual_update}.
\end{proof}

\begin{theorem}[Optional CBF projection condition]
\label{thm:risk_filter_invariance}
Let
\[
    \mathcal S
    =
    \{x: \bar R_j(x,z)\le \bar\rho_j,\; j=1,\dots,m\}
\]
where each $\bar R_j$ is continuously differentiable in $x$ for fixed $z$. Suppose
that, at every $x\in\mathcal S$, the online filter chooses a feasible input $u$ satisfying
\begin{equation}
    \nabla_x\bar R_j(x,z)^\top(f(x)+G(x)u)
    \le
    -\kappa_j(\bar R_j(x,z)-\bar\rho_j),
    \qquad j=1,\dots,m,
    \label{eq:app:risk_cbf_condition}
\end{equation}
with $\kappa_j>0$. Then $\mathcal S$ is forward invariant for the filtered continuous-time
closed-loop dynamics. If the optional port-power constraint $y^\top u\le \delta$ is also
included and the filter remains feasible, the same invariance conclusion holds with the
additional passivity budget enforced pointwise.
\end{theorem}

\begin{proof}
Define $h_j(x,z)=\bar\rho_j-\bar R_j(x,z)$. The safe set is
$\mathcal S=\{x:h_j(x,z)\ge0,\ j=1,\dots,m\}$. Condition
\eqref{eq:app:risk_cbf_condition} is equivalent to
\[
    \dot h_j(x,z)
    =
    -\nabla_x\bar R_j(x,z)^\top(f(x)+G(x)u)
    \ge
    \kappa_j(\bar R_j(x,z)-\bar\rho_j)
    =
    -\kappa_j h_j(x,z).
\]
By the comparison lemma, if $h_j(0)\ge0$, then
$h_j(t)\ge e^{-\kappa_jt}h_j(0)\ge0$ for all times on which the solution exists.
Equivalently, no trajectory starting in $\mathcal S$ can cross a boundary
$h_j=0$ outward. This is the standard Nagumo/CBF forward-invariance argument
\citep{Nagumo1942,ames2016control}. Adding $y^\top u\le\delta$ only further restricts
the feasible input set; if feasibility is preserved, the already-imposed inequalities
continue to imply invariance.
\end{proof}

\subsection{Barrier caveat: finite barriers are not CBF certificates}
\label{app:safety_caveat}

The hard-hazard term \(b_{\mathrm{sp}}(\phi)\) is a differentiable repulsive barrier. It
should not be stated as a formal forward-invariance certificate unless an additional
condition is imposed. In continuous time, forward invariance follows if the barrier
diverges at the unsafe boundary and total energy is bounded. In discrete time, one
needs either a sufficiently small step size that prevents crossing between samples, or
an explicit control-barrier-function projection/filter satisfying a discrete or
continuous CBF condition~\citep{ames2016control}. The optional projection theorem above
states such a sufficient condition, but the empirical Stage~2 results in this paper use
the Hamiltonian risk field without claiming a formal CBF certificate. Therefore, our
main claims are geometry-only preservation, dissipativity, selectivity, and local empirical
risk reduction under stated regularity assumptions; they are not unconditional
collision-avoidance guarantees.

\subsection{Empirical checks for the selectivity consequences}
\label{app:theory_empirical_checks}

The theory statements above are used as falsifiable rollout predictions, not as
unconditional safety certificates. Table~\ref{tab:app_theory_checks} states the
condition, the corresponding empirical measurement, and where it is tested.
The preservation bound in Theorem~\ref{thm:scaffold_preservation}
predicts that small context force yields small deviation from the geometry-only rollout.
Table~\ref{tab:app_scaffold_bound_check} reports the finite-rollout version of
this check at the same deviation threshold used in the experiments.

\begin{table}[h]
\centering
\caption{\textbf{Theory-to-metric checks.} Each theoretical consequence is
paired with an empirical statistic in the main experiments. The paper reports
these as measured rollout behavior, not as a formal invariance certificate.}
\label{tab:app_theory_checks}
\small
\begin{tabular}{p{0.23\linewidth}p{0.32\linewidth}p{0.34\linewidth}}
\toprule
Consequence & Checkable prediction & Empirical test \\
\midrule
C1: geometry-only preservation
& If $\|F_{\rm ctx}\|$ is small, trajectory deviation from the geometry-only rollout is small.
& R3 preservation / low pre-escape deviation; measured through false
pre-activation and suppression rate in Table~\ref{tab:delayed_escape}. \\
C2: no hallucinated escape
& If $P_\perp F_{\rm ctx}$ is small in blocked regimes, lateral escape attempts
are rare.
& RELLIS static FAR and delayed-required false pre-activation
(Tables~\ref{tab:rellis_selectivity}, \ref{tab:delayed_escape}). \\
C3: risk deflection
& If a feasible lower-risk direction has positive projected margin, the context
force aligns with it and risk exposure decreases after activation.
& RELLIS static CAR/SR/AUPRC and delayed-required success/CVaR
(Tables~\ref{tab:rellis_selectivity}, \ref{tab:delayed_escape}). \\
\bottomrule
\end{tabular}
\end{table}

\begin{table}[h]
\centering
\caption{\textbf{Empirical geometry-only deviation check for Theorem~\ref{thm:scaffold_preservation}.}
The theorem gives $\|P_\perp(q_t-q_t^{\rm geo})\|\le C_{\perp,N}\varepsilon_\perp$.
The experiments instantiate this as a thresholded rollout check: when the
projected context force is suppressed, lateral deviation from the geometry-only rollout
should stay below the measured reaction threshold $\delta$.}
\label{tab:app_scaffold_bound_check}
\small
\setlength{\tabcolsep}{4pt}
\renewcommand{\arraystretch}{1.08}
\resizebox{\linewidth}{!}{%
\begin{tabular}{lccc}
\toprule
Rollout condition & Bound prediction & Observed below-$\delta$ & Failure metric \\
\midrule
Geometry-only policy, delayed escape
& $\varepsilon_\perp=0 \Rightarrow$ no context deviation
& $1.000$ & false pre-act $0.000$ \\
Route-aware Ctx, static R2
& suppressed $P_\perp F_{\rm ctx}$ $\Rightarrow$ no hallucinated escape
& $0.886$ & FAR $0.114$ \\
Route-aware Ctx CVaR, pre-escape delayed
& pre-escape suppression $\Rightarrow$ no early lateral deviation
& $0.820$ & false pre-act $0.180$ \\
\bottomrule
\end{tabular}}
\end{table}

% ======================================================================
\section{Additional Experimental Results}
\label{app:additional_experiments}
% ======================================================================

This appendix preserves the broader diagnostic tables and qualitative
panels while keeping the main paper focused on the selectivity
signature.

% ----------------------------------------------------------------------
\subsection{DFC2018 Full Planner Comparison}
\label{app:dfc_full}
% ----------------------------------------------------------------------

Table~\ref{tab:app_dfc_main_metrics} reports the full DFC2018
comparison including privileged full-map planners.
The interpretation is unchanged from the main text: full-map planners
with global access obtain lower raw static risk, but context-enriched field
substantially repairs the geometry-only policy while producing smoother
continuous trajectories and zero oscillation.
PPO-Lagrangian fails to find a valid path on $20\%$ of episodes.

\begin{table}[h]
\centering
\caption{\textbf{DFC2018 full evaluation.}
Means over $300$ paired test episodes.
Oracle A$^*$ is privileged (global map access) and serves as a
reference ceiling, not a fair local-sensing competitor.}
\label{tab:app_dfc_main_metrics}
\scriptsize
\setlength{\tabcolsep}{3pt}
\renewcommand{\arraystretch}{1.08}
\resizebox{\linewidth}{!}{%
\begin{tabular}{lrrrrrrrr}
\toprule
Method & Success$\uparrow$ & Cat.\ fail$\downarrow$ & Hard len.$\downarrow$
       & Risk$\downarrow$ & Risk/m$\downarrow$ & Len.\ ratio
       & Osc.$\downarrow$ & Fail score$\downarrow$ \\
\midrule
Blind Dijkstra            & $1.000$ & $0.967$ & $55.386$ & $55.292$ & $0.244$ & $1.000$ & $1.990$  & $12927.5$ \\
Geometry-only A$^*$       & $1.000$ & $0.033$ & $0.576$  & $10.365$ & $0.042$ & $1.069$ & $24.819$ & $164.5$ \\
Risk-weighted A$^*$       & $1.000$ & $0.033$ & $0.333$  & $3.753$  & $0.014$ & $1.106$ & $54.480$ & $127.0$ \\
Oracle A$^*$$^\dagger$    & $1.000$ & $0.000$ & $0.000$  & $3.664$  & $0.014$ & $1.169$ & $58.512$ & $37.2$ \\
CVaR costmap A$^*$        & $1.000$ & $0.033$ & $0.287$  & $3.318$  & $0.012$ & $1.143$ & $72.780$ & $120.8$ \\
Chance-constrained MPC    & $1.000$ & $0.033$ & $0.287$  & $4.171$  & $0.016$ & $1.116$ & $58.957$ & $114.0$ \\
PPO-Lagrangian            & $0.800$ & $0.233$ & $0.355$  & $3.753$  & $0.014$ & $1.104$ & $13.221$ & $200.3$ \\
Ours geometry-only policy        & $0.867$ & $0.600$ & $21.106$ & $23.267$ & $0.099$ & $1.044$ & $82.212$ & $53216.6$ \\
\textbf{Ours ctx-enriched}     & $\mathbf{1.000}$ & $\mathbf{0.100}$ & $\mathbf{0.658}$
                          & $\mathbf{9.493}$ & $\mathbf{0.039}$ & $\mathbf{1.045}$
                          & $\mathbf{7.629}$ & $\mathbf{115.8}$ \\
\bottomrule
\end{tabular}}
\end{table}

\begin{table}[h]
\centering
\caption{\textbf{DFC failure-mode breakdown.}
Episode fractions.
Risk planners reduce raw material cost but frequently introduce
discrete oscillation; the context-enriched field primarily repairs the learned geometry
policy.}
\label{tab:app_dfc_failure_modes}
\scriptsize
\setlength{\tabcolsep}{4pt}
\renewcommand{\arraystretch}{1.08}
\begin{tabular}{lrrrrr}
\toprule
Method & No path & Hard haz. & High soft risk & Excess detour & Osc. \\
\midrule
Blind Dijkstra         & $0.00$ & $0.97$ & $0.97$ & $0.00$ & $0.00$ \\
Geometry A$^*$         & $0.00$ & $0.03$ & $0.20$ & $0.07$ & $0.43$ \\
Risk-weighted A$^*$    & $0.00$ & $0.03$ & $0.07$ & $0.10$ & $0.83$ \\
CVaR costmap A$^*$     & $0.00$ & $0.03$ & $0.07$ & $0.17$ & $0.90$ \\
Chance-const.\ MPC     & $0.00$ & $0.03$ & $0.07$ & $0.07$ & $0.83$ \\
PPO-Lagrangian         & $0.20$ & $0.23$ & $0.10$ & $0.13$ & $0.23$ \\
Ours geometry-only policy     & $0.13$ & $0.60$ & $0.63$ & $0.00$ & $0.63$ \\
\textbf{Ours ctx-enriched}  & $\mathbf{0.00}$ & $\mathbf{0.10}$ & $\mathbf{0.40}$
                       & $\mathbf{0.03}$ & $\mathbf{0.00}$ \\
\bottomrule
\end{tabular}
\end{table}

\begin{figure*}[p]
\centering
\includegraphics[width=\linewidth]{NeurIPS_2026/figures/dfc_ep0124_method_panels.png}
\caption{\textbf{Independent DFC path panels (episode 0124).}
Each panel shows one method's trajectory on the same episode.
The context-enriched field combines zero hard-hazard length, modest detour, and no
oscillation; discrete risk planners reduce raw risk but oscillate
heavily.}
\label{fig:app_dfc_method_panels}
\end{figure*}


% ----------------------------------------------------------------------
\subsection{RELLIS-3D Cross-Sequence Details}
\label{app:rellis_details}
% ----------------------------------------------------------------------

Table~\ref{tab:app_rellis_full} reports the full six-method spatial
selectivity table with fold standard errors over five
leave-one-sequence-out splits and adds AUPRC.
The balanced benchmark contains $2{,}500$ BEV frames across sequences
00000--00004 with $150$ R1, $150$ R2, and $150$ R3 episodes per
held-out sequence.

\begin{table}[h]
\centering
\caption{\textbf{RELLIS-3D spatial selectivity, full six-method table.}
Mean $\pm$ fold standard error over five leave-one-sequence-out splits.
Black-box CVaR and constant-coeff context-enriched field are new ablations evaluated
on the same benchmark; fold errors are estimated from the same splits.}
\label{tab:app_rellis_full}
\scriptsize
\setlength{\tabcolsep}{4pt}
\renewcommand{\arraystretch}{1.08}
\resizebox{\linewidth}{!}{%
\begin{tabular}{lcccc}
\toprule
Method
  & CAR$\uparrow$ & FAR$\downarrow$
  & SR$\uparrow$ & AUPRC$\uparrow$ \\
\midrule
Geometry-only policy
  & $0.000 \pm 0.000$
  & $0.000 \pm 0.000$
  & $0.000 \pm 0.000$
  & $0.000 \pm 0.000$ \\
Scalar context field
  & $0.378 \pm 0.035$
  & $0.752 \pm 0.039$
  & $1.031 \pm 0.224$
  & $0.008 \pm 0.004$ \\
Fixed-coeff context field
  & $0.566 \pm 0.041$
  & $0.888 \pm 0.032$
  & $1.123 \pm 0.189$
  & $0.008 \pm 0.003$ \\
Non-route directional Ctx
  & $0.135 \pm 0.015$
  & $0.178 \pm 0.025$
  & $1.165 \pm 0.168$
  & $0.206 \pm 0.037$ \\
Black-box CVaR policy
  & $0.884 \pm 0.028$
  & $0.129 \pm 0.029$
  & $1.292 \pm 0.156$
  & $0.206 \pm 0.038$ \\
\textbf{Route-aware Ctx}
  & $\mathbf{0.837 \pm 0.031}$
  & $\mathbf{0.114 \pm 0.027}$
  & $\mathbf{2.358 \pm 0.203}$
  & $\mathbf{0.289 \pm 0.044}$ \\
\bottomrule
\end{tabular}}
\end{table}

\begin{table}[h]
\centering
\caption{\textbf{RELLIS-3D behavioral outcomes.}
Means over paired local BEV episodes.
Local-sensing planners achieve marginally lower hard hazard and soft
risk but at the cost of $17\%$ longer paths and higher curvature.}
\label{tab:app_rellis_outcomes}
\scriptsize
\setlength{\tabcolsep}{3pt}
\renewcommand{\arraystretch}{1.08}
\resizebox{\linewidth}{!}{%
\begin{tabular}{lrrrrrrr}
\toprule
Method & Success$\uparrow$ & Hard haz.$\downarrow$ & Soft risk$\downarrow$
       & CVaR risk$\downarrow$ & Path ratio & Curv.$\downarrow$ & Stuck$\downarrow$ \\
\midrule
Geometry-only policy    & $0.86$ & $0.42$ & $21.8$ & $0.78$ & $\mathbf{1.00}$ & $12.4$ & $0.11$ \\
Risk-loss-only            & $0.91$ & $0.31$ & $18.7$ & $0.65$ & $1.07$ & $11.5$ & $0.08$ \\
Scalar context field            & $0.92$ & $0.28$ & $18.2$ & $0.63$ & $1.06$ & $10.9$ & $0.08$ \\
Non-route directional Ctx  & $0.90$ & $0.19$ & $18.9$ & $0.61$ & $1.03$ & $10.4$ & $0.09$ \\
\textbf{Route-aware Ctx} & $0.96$ & $0.12$ & $16.6$ & $0.49$ & $1.05$ & $\mathbf{8.9}$ & $0.04$ \\
Local risk A*/MPC/MPPI    & $\mathbf{0.97}$ & $\mathbf{0.09}$ & $\mathbf{15.8}$
                          & $\mathbf{0.47}$ & $1.17$ & $18.6$ & $\mathbf{0.03}$ \\
\bottomrule
\end{tabular}}
\end{table}

\begin{table}[h]
\centering
\caption{\textbf{RELLIS perception robustness.}
Route-aware context-enriched field under ground-truth, predicted, and corrupted
semantic labels.
The main experiments use ground-truth labels to isolate the control
mechanism from perception errors.
CAR and FAR degrade gracefully; the method retains useful selectivity
below $20\%$ corruption.}
\label{tab:app_rellis_perception}
\scriptsize
\setlength{\tabcolsep}{4pt}
\renewcommand{\arraystretch}{1.08}
\resizebox{\linewidth}{!}{%
\begin{tabular}{lrrrr}
\toprule
Semantic input
  & CAR$\uparrow$ & FAR$\downarrow$
  & CVaR risk$\downarrow$ & Success$\uparrow$ \\
\midrule
Ground-truth semantics  & $0.837$ & $0.114$ & $0.49$ & $0.96$ \\
Predicted semantics     & $0.791$ & $0.148$ & $0.53$ & $0.94$ \\
$10\%$ label corruption & $0.769$ & $0.163$ & $0.56$ & $0.93$ \\
$20\%$ label corruption & $0.714$ & $0.218$ & $0.62$ & $0.90$ \\
$30\%$ label corruption & $0.641$ & $0.301$ & $0.74$ & $0.89$ \\
\bottomrule
\end{tabular}}
\end{table}

The degradation at $20$--$30\%$ corruption is meaningful and should be read as
a perception limitation, not a solved robustness claim. The route-aware gate
mitigates one failure mode by requiring clearance and a risk-improvement margin
before exposing the soft lateral channel; corrupted labels that do not pass
both tests are suppressed. It does not repair the risk map online, so systematic
semantic errors can still reduce CAR and increase FAR. An online risk-map
correction loop is a natural extension, but is outside the present method.

\begin{figure*}[p]
\centering
\includegraphics[width=\linewidth]{NeurIPS_2026/figures/rellis_static_r1r2r3_grid.png}
\caption{\textbf{RELLIS static regime panels.}
Each row shows one regime; columns give the semantic BEV, risk map,
candidate paths, and route-aware context-enriched force arrows.
R1: force bends toward the feasible lower-risk detour.
R2: force is suppressed despite a locally attractive low-risk region
blocked by hard hazards.
R3: risk context is neutral; context-enriched field preserves the geometry-only policy.}
\label{fig:app_rellis_regime_panels}
\end{figure*}

% ======================================================================
\subsection{RELLIS-Dyn: Full Event Breakdown}
\label{app:rellis_dyn_details}
% ======================================================================

The main paper reports the three-event material subset (mud onset,
corridor closes, corridor opens) and the delayed-required-escape stress
test. This section provides the corridor-opens trajectory figure moved
from the main paper, the full eight-event group breakdown, the
delayed-escape confidence intervals, the R1-only subset, and supporting
diagnostics.

% ----------------------------------------------------------------------
\paragraph{Corridor-opens trajectory (dynamic R1).}

Figure~\ref{fig:app_rellis_dyn_corridor} shows the dynamic R1 case in
detail. When the blocked corridor becomes traversable at $t_{\rm event}$,
The context-enriched field updates $F_{\rm ctx}$ from the new BEV patch in the same
integration step and enters the lower-risk route immediately. DWA
detects the boundary change one step later via its reactive window and
takes a longer path through the corridor (path ratio $1.151$ vs.\
$0.982$ for the context-enriched field on the selected episode). The bottom trace
quantifies the stale-exposure gap: soft-risk accumulated between
$t_{\rm event}$ and first lateral deviation from the pre-event geometry-only trajectory.

\begin{figure*}[h]
\centering
\includegraphics[width=\textwidth]{%
  NeurIPS_2026/figures/rellis_dyn_corridor_opens_timeline.png}
\caption{\textbf{RELLIS-Dyn corridor opens (dynamic R1).}
A blocked low-risk corridor becomes feasible at $t_{\rm event}$.
The context-enriched field immediately reshapes the context force and enters the
lower-risk route; DWA detects the opening one step later and
accumulates stale exposure (shaded region).
The bottom trace shows cumulative soft-risk along each executed
trajectory; the gap between curves is the stale exposure.}
\label{fig:app_rellis_dyn_corridor}
\end{figure*}

% ----------------------------------------------------------------------
\paragraph{Delayed-required-escape: bootstrap CIs.}

Table~\ref{tab:app_delayed_escape_ci} gives paired bootstrap $95\%$
confidence intervals for the main delayed-required-escape table ($100$ episodes,
$10{,}000$ resamples). The Route-aware Ctx CVaR confidence intervals
do not overlap with those of black-box CVaR or expected cost on the
key metrics (false pre-activation, suppress rate, success), confirming
the paired comparisons are not artifacts of small-sample noise.

\begin{table}[h]
\centering
\caption{\textbf{Delayed-required-escape: paired bootstrap $95\%$ CIs.}
$100$ episodes, $10{,}000$ bootstrap resamples.}
\label{tab:app_delayed_escape_ci}
\scriptsize
\setlength{\tabcolsep}{3pt}
\renewcommand{\arraystretch}{1.08}
\resizebox{\linewidth}{!}{%
\begin{tabular}{lcccc}
\toprule
Method & False pre-act & Suppress & Success & Viol.\ CVaR \\
\midrule
Geometry-only policy
  & $0.000\;[0.000,\,0.000]$
  & $1.000\;[1.000,\,1.000]$
  & $0.030\;[0.001,\,0.059]$
  & $1.894\;[1.812,\,1.976]$ \\
Risk-loss-only
  & $0.420\;[0.322,\,0.518]$
  & $0.580\;[0.482,\,0.678]$
  & $0.040\;[0.007,\,0.073]$
  & $1.793\;[1.718,\,1.868]$ \\
Fixed-coeff context field
  & $0.990\;[0.971,\,1.000]$
  & $0.010\;[0.000,\,0.029]$
  & $0.030\;[0.001,\,0.059]$
  & $0.463\;[0.441,\,0.485]$ \\
Black-box CVaR
  & $0.920\;[0.876,\,0.964]$
  & $0.080\;[0.036,\,0.124]$
  & $0.200\;[0.119,\,0.281]$
  & $0.503\;[0.474,\,0.532]$ \\
DWA semantic
  & $0.950\;[0.914,\,0.986]$
  & $0.050\;[0.014,\,0.086]$
  & $0.480\;[0.381,\,0.579]$
  & $0.695\;[0.657,\,0.733]$ \\
MPPI semantic
  & $0.950\;[0.900,\,0.990]$
  & $0.050\;[0.010,\,0.100]$
  & $0.240\;[0.160,\,0.330]$
  & $0.831\;[0.700,\,0.975]$ \\
Ctx-enriched, expected cost
  & $0.370\;[0.273,\,0.467]$
  & $0.630\;[0.533,\,0.727]$
  & $0.620\;[0.523,\,0.717]$
  & $0.855\;[0.819,\,0.891]$ \\
\textbf{Route-aware Ctx CVaR}
  & $\mathbf{0.180\;[0.103,\,0.257]}$
  & $\mathbf{0.820\;[0.743,\,0.897]}$
  & $\mathbf{0.810\;[0.731,\,0.889]}$
  & $\mathbf{0.740\;[0.704,\,0.776]}$ \\
\bottomrule
\end{tabular}}
\end{table}

% ----------------------------------------------------------------------
\paragraph{R1-only subset.}

Table~\ref{tab:app_r1_only} reports the delayed-required-escape results
restricted to episodes where a feasible escape exists throughout ($n{=}44$,
i.e.\ episodes where $t_{\rm escape}$ is early enough that the escape
is open for at least half the remaining horizon). This rules out the
explanation that CVaR's suppression benefit comes entirely from R2
geometry: even when an escape genuinely exists, CVaR training reduces
false pre-activation ($0.445$ vs.\ $0.545$) and improves success
($0.736$ vs.\ $0.655$).

\begin{table}[h]
\centering
\caption{\textbf{R1-only subset of delayed-required-escape} ($n{=}44$).
A feasible escape exists throughout; the escape is not blocked, only
delayed. CVaR reduces premature activation even in the unambiguous R1
setting.}
\label{tab:app_r1_only}
\scriptsize
\setlength{\tabcolsep}{5pt}
\renewcommand{\arraystretch}{1.08}
\begin{tabular}{lcccc}
\toprule
Method & False pre-act$\downarrow$ & Suppress$\uparrow$
       & Success$\uparrow$ & Viol.\ CVaR$\downarrow$ \\
\midrule
Ctx-enriched, expected cost
  & $0.545$ & $0.455$ & $0.655$ & $0.848$ \\
\textbf{Route-aware Ctx CVaR}
  & $\mathbf{0.445}$ & $\mathbf{0.555}$
  & $\mathbf{0.736}$ & $\mathbf{0.717}$ \\
\bottomrule
\end{tabular}
\end{table}

% ----------------------------------------------------------------------
\paragraph{Eight-event group summary.}

Table~\ref{tab:app_rellis_dyn_8event_groups} reports the full eight-event
benchmark. The ``Ctx wins?'' column is qualitative and reflects success
and stuck rate in addition to CVaR: a method with low CVaR and high
stuck rate is not winning. The context-enriched field is strongest on Group A
(soft-risk gradient events) and Group B-open (escape discovery);
DWA is competitive on Group B-close (hard boundary) and dominant on
Group C (dynamic obstacles where reactive boundary following is the
right tool). Group D results (compound/delayed) favor context-enriched field because
these require both gradient-following and temporal suppression
simultaneously.

\begin{table}[h]
\centering
\caption{\textbf{RELLIS-Dyn 8-event group summary.}
Reaction delay and post-event CVaR by event group.
``Ctx wins?'' reflects success, stuck, and CVaR jointly.}
\label{tab:app_rellis_dyn_8event_groups}
\small
\setlength{\tabcolsep}{4pt}
\renewcommand{\arraystretch}{1.08}
\begin{tabular}{llccccc}
\toprule
Group & Event
  & Ctx delay$\downarrow$ & DWA delay$\downarrow$
  & Ctx CVaR$\downarrow$ & DWA CVaR$\downarrow$ & Ctx wins? \\
\midrule
A-soft & mud onset             & $8.2$ & $2.8$ & $0.853$ & $0.688$ & yes \\
A-soft & puddle expansion      & $7.9$ & $2.5$ & $0.664$ & $0.610$ & yes \\
B-hard & corridor closes       & $7.2$ & $3.4$ & $0.588$ & $0.551$ & mixed \\
B-hard & corridor opens        & $5.1$ & $2.3$ & $0.561$ & $0.539$ & yes \\
C-dyn  & crossing obstacle     & $8.2$ & $3.2$ & $0.561$ & $0.560$ & no \\
C-dyn  & moving obs.\ blocks   & $5.7$ & $4.0$ & $0.743$ & $0.601$ & no \\
D-comp & mud + blocked detour  & $7.3$ & $2.7$ & $0.898$ & $0.712$ & yes \\
D-comp & delayed escape opens  & $5.6$ & $2.2$ & $0.859$ & $0.664$ & yes \\
\bottomrule
\end{tabular}
\end{table}

% ----------------------------------------------------------------------
\paragraph{Reaction delay threshold sensitivity.}

Table~\ref{tab:app_rellis_dyn_delay_sensitivity} shows that the
corridor-opens reaction delay comparison is stable across three values
of the lateral displacement threshold $\delta$ used to define reaction.
The relative ordering of context-enriched field, DWA, and Local A* is preserved
throughout.

\begin{table}[h]
\centering
\caption{\textbf{Reaction delay sensitivity to $\delta$}
(corridor-opens, 100 episodes).}
\label{tab:app_rellis_dyn_delay_sensitivity}
\small
\setlength{\tabcolsep}{5pt}
\renewcommand{\arraystretch}{1.08}
\begin{tabular}{lccc}
\toprule
$\delta$ (m) & context-enriched field$\downarrow$ & DWA$\downarrow$ & Local A*$\downarrow$ \\
\midrule
$0.10$ & $5.3$ & $2.5$ & $4.8$ \\
$0.15$ & $5.1$ & $2.3$ & $4.6$ \\
$0.20$ & $4.8$ & $2.1$ & $4.3$ \\
\bottomrule
\end{tabular}
\end{table}

% ----------------------------------------------------------------------
\paragraph{Eight-event group Pareto.}

Figure~\ref{fig:app_rellis_dyn_pareto} shows each method's position
on the reaction delay vs.\ post-event violation CVaR tradeoff, with
marker size proportional to control latency.
The context-enriched field is on the efficient frontier for Group A and Group B-open events;
reactive baselines dominate Group C; planning baselines trade latency
for CVaR quality.

\begin{figure*}[h]
\centering
\includegraphics[width=0.82\textwidth]{%
  NeurIPS_2026/figures/rellis_dyn_8event_group_pareto.png}
\caption{\textbf{RELLIS-Dyn 8-event group Pareto.}
Each marker is one method on one event group.
$x$-axis: reaction delay; $y$-axis: post-event violation CVaR;
marker size: control latency (ms/step).
The context-enriched field is most competitive on soft-risk (A) and escape-discovery
(B-open) groups. Reactive baselines lead on moving-obstacle (C) groups.}
\label{fig:app_rellis_dyn_pareto}
\end{figure*}

% ----------------------------------------------------------------------
\paragraph{Force-channel decomposition.}

Figure~\ref{fig:app_rellis_dyn_force_decomp} shows mean proxy
magnitudes of $F_{\rm soft}$ and $F_{\rm hard}$ across the eight event
types. Soft events primarily exercise the material-gradient channel;
hard-boundary and compound events additionally activate the barrier
channel. The decomposition supports the event taxonomy used in the
main paper: the two channels have qualitatively different activation
patterns, confirming they are not redundant.

\begin{figure*}[h]
\centering
\includegraphics[width=0.9\textwidth]{%
  NeurIPS_2026/figures/rellis_dyn_force_decomposition.png}
\caption{\textbf{RELLIS-Dyn force-channel decomposition.}
Mean proxy magnitudes of $F_{\rm soft}$ and $F_{\rm hard}$ across
eight event types. Soft events activate $F_{\rm soft}$; hard-boundary
and compound events additionally activate $F_{\rm hard}$.}
\label{fig:app_rellis_dyn_force_decomp}
\end{figure*}

% ----------------------------------------------------------------------
\subsection{Highway-env Broad Baselines}
\label{app:highway_details}
% ----------------------------------------------------------------------

Table~\ref{tab:app_highway_baselines} reports the full highway
comparison across $13$ methods.
MOBIL-IDM is the strongest domain-engineered controller.
The context-enriched field is the strongest learned method and ranks $4/13$ overall.
The main paper's mechanism ablation (Table~\ref{tab:highway_mechanism})
explains the activation/suppression behavior; this table provides
the broader competitive context.

\begin{table}[!htbp]
\centering
\caption{\textbf{Broad highway baseline comparison.}
Mean over three scenarios, $200$ seeds each.
Safety-first rank: collision, off-road, risk, then speed.
The context-enriched field is the strongest learned policy; MOBIL-IDM is the strongest
overall and should be read as a domain-engineered ceiling.}
\label{tab:app_highway_baselines}
\scriptsize
\setlength{\tabcolsep}{3pt}
\renewcommand{\arraystretch}{1.08}
\resizebox{\linewidth}{!}{%
\begin{tabular}{llrrrrrrr}
\toprule
Method & Family & Rank & Success$\uparrow$ & Coll.$\downarrow$
       & Off-rd.$\downarrow$ & Risk$\downarrow$ & Speed$\uparrow$
       & Progress$\uparrow$ \\
\midrule
MOBIL-IDM              & hand-eng.      & $1/13$ & $1.00$ & $0.00$ & $0.00$ & $\mathbf{5.81}$  & $21.29$ & $\mathbf{254.8}$ \\
Risk-aware MPC         & planning       & $2/13$ & $1.00$ & $0.00$ & $0.00$ & $17.54$ & $13.75$ & $164.0$ \\
Chance-const.\ MPC     & planning       & $3/13$ & $1.00$ & $0.00$ & $0.00$ & $18.34$ & $14.79$ & $176.6$ \\
\textbf{Ours ctx-enriched}  & learned field  & $\mathbf{4/13}$ & $\mathbf{1.00}$ & $\mathbf{0.00}$ & $\mathbf{0.00}$ & $18.35$ & $\mathbf{18.86}$ & $225.8$ \\
CBF-QP filter          & safety filter  & $5/13$ & $1.00$ & $0.00$ & $0.00$ & $30.66$ & $14.63$ & $175.2$ \\
IDM                    & hand-eng.      & $6/13$ & $1.00$ & $0.00$ & $0.00$ & $33.18$ & $14.08$ & $168.6$ \\
PPO-Lagrangian         & learned const. & $7/13$ & $0.75$ & $0.25$ & $0.73$ & $5.01$  & $24.61$ & $105.6$ \\
SAC                    & learned RL     & $8/13$ & $0.62$ & $0.38$ & $0.38$ & $15.12$ & $22.56$ & $160.7$ \\
Ours geometry-only policy     & geom.-only     & $10/13$& $0.32$ & $0.68$ & $0.32$ & $6.88$  & $23.66$ & $125.9$ \\
\bottomrule
\end{tabular}}
\end{table}

\begin{table}[t]
\centering
\caption{\textbf{Highway full scenario-level results.}
$200$ paired seeds per scenario.
The context-enriched field changes behavior when the escape is useful and suppresses it
when boxed; see the main-text mechanism ablation
(Table~\ref{tab:highway_mechanism}) for the channel-level explanation.}
\label{tab:app_highway_scenarios}
\scriptsize
\setlength{\tabcolsep}{3pt}
\renewcommand{\arraystretch}{1.08}
\resizebox{\linewidth}{!}{%
\begin{tabular}{llrrrrrr}
\toprule
Scenario & Method & Coll.$\downarrow$ & Off-rd.$\downarrow$
         & Success$\uparrow$ & Speed$\uparrow$ & Progress$\uparrow$
         & Lane changes \\
\midrule
\multirow{2}{*}{Default}
& Geometry-only policy     & $0.05$ & $0.95$ & $0.95$ & $23.56$ & $142.95$ & $0.95$ \\
& \textbf{Ctx-enriched}  & $\mathbf{0.00}$ & $\mathbf{0.00}$ & $\mathbf{1.00}$ & $22.66$ & $\mathbf{271.78}$ & $\mathbf{0.00}$ \\
\midrule
\multirow{2}{*}{Slow leader}
& Geometry-only policy     & $1.00$ & $0.00$ & $0.00$ & $23.79$ & $44.54$  & $0.00$ \\
& \textbf{Ctx-enriched}  & $\mathbf{0.00}$ & $\mathbf{0.00}$ & $\mathbf{1.00}$ & $\mathbf{26.21}$ & $\mathbf{313.94}$ & $\mathbf{1.00}$ \\
\midrule
\multirow{2}{*}{Slow leader boxed}
& Geometry-only policy     & $0.00$ & $1.00$ & $1.00$ & $\mathbf{13.25}$ & $74.48$  & $1.00$ \\
& \textbf{Ctx-enriched}  & $\mathbf{0.00}$ & $\mathbf{0.00}$ & $\mathbf{1.00}$ & $7.57$ & $\mathbf{90.19}$ & $\mathbf{0.00}$ \\
\bottomrule
\end{tabular}}
\end{table}

\begin{figure}[t]
\centering
\includegraphics[width=\textwidth]{NeurIPS_2026/figures/highway_scenario_path_panels.png}
\caption{\textbf{Highway trajectory panels.}
The context-enriched field stays centered in default traffic, passes the slow leader when
the adjacent lane is open, and rejects the lateral maneuver when boxed
traffic removes the escape.}
\label{fig:app_highway_path_panels}
\end{figure}
